% ---------------------------------------------------------------------------
% Chapter 1: The Ancestral Home: Protons, PDFs and the Hard Scattering
% The scripts, configurations and records discussed here live in example/ME/
% of https://github.com/ravindkv/JitJet . The book gives the physics and the
% instructions; the code itself is read on GitHub, not printed here.
% ---------------------------------------------------------------------------
\chapter{The Ancestral Home: Protons, PDFs and the Hard Scattering}
\chaptermark{The Ancestral Home}
\label{ch:ancestral_home}

\begin{journeybox}
Every family has an ancestral home. Ours is the proton, and the mother of our
jet, a bottom quark, is born there in a hard scattering. Before she takes a
single step we must know how much energy she was given, because everything
that follows on this journey is measured against it.
\end{journeybox}

\physicsline{factorisation, PDFs and their DGLAP evolution, the leading-order
matrix element for $q\bar q\to\bbbar$, the change of variables behind every
Monte Carlo integrator, the running coupling and the scale choice, the
seven-point scale variation, \PYTHIA's \code{HardQCD:qqbar2bbbar} process
inside \CMSSW, the same process in \MG, the hard-process record of our event}

\section{Welcome aboard: what we will see at this first stop}
\label{sec:ancestral_home:what_we_compute}

Welcome, and please take a seat. I am your guide for this journey, and I have
one request before we set off: do not just look out of the window. At every
stop I will hand you the formulae, the numbers and the buttons, and I would
like you to press the buttons yourself. A jet is not something you can
understand by being told about it; it is something you follow.

Our first stop is the ancestral home, and there is exactly one number to take
away from it: the cross section for producing a $\bbbar$ pair with
$\pthat\ge100\GeV$ in proton--proton collisions at $\sqrt{s}=13.6\TeV$
through quark--antiquark annihilation. It comes out at about $345\pb$. We
will meet it \emph{four} times, and the order in which we meet it is the
order of every stop in this book, from the top down: first you watch, then
you run, then you derive, and only then do you check against the official
machinery.
\begin{enumerate}[nosep]
  \item \textbf{Watch it.} A short film, rendered from the same code that
    prints the cross section, shows the collision, the integration volume
    and the number converging (\cref{sec:ancestral_home:animation}).
  \item \textbf{Run it.} A single Python file, needing only NumPy and SciPy,
    does the integral in about one second; you press the button and watch
    $345\pb$ come out on your own screen
    (\cref{sec:ancestral_home:standalone_code}).
  \item \textbf{Derive it.} From the factorisation formula, the parton
    distribution functions and a one-line matrix element we write down the
    three-dimensional integral the script evaluates and understand every
    factor in it
    (\crefrange{sec:ancestral_home:two_protons_meet}{sec:ancestral_home:from_formula_to_number}).
  \item \textbf{Check it against the official packages.} The CMS generator
    configuration runs the same process inside \PYTHIA and its log prints
    the cross section it sampled; \MG, a second and independent generator,
    computes the same Born process from its own Feynman diagrams
    (\cref{sec:ancestral_home:validation}).
\end{enumerate}
\Cref{tab:ancestral_home:three_numbers} gives away the ending. The four agree
within the statistical precision of the Monte Carlo runs, and that agreement
is the whole point of the stop: after this chapter you will know what every
knob in the CMS configuration does to the number, because you will have turned
each knob yourself.

\paragraph{Top to bottom, bottom to top.} A word on the order, because it
is the opposite of a textbook's. A textbook derives first and shows last; a
guided tour shows first. So at every stop you will see the film before the
formula, run the script before you can read it, and only then sit down with
pencil and paper to earn what you have already seen, before the official
packages get the last word. The journey itself runs the other way, from the
bottom up: the ``bottom to top'' of the title is the jet's whole road from
the hard scattering to the histogram, and it is also the two quarks we
follow, the bottom quark of this chapter and the top quark of
\cref{ch:three_families}. The two directions do not conflict. The road is
walked from its beginning to its end; each stop along it is explained from
the top down.

\paragraph{How each derivation is told.} The theory part of every stop is
itself told the same way, section after section, and it is worth knowing the
pattern before you meet it. Each section opens with a question you should
try to answer before reading on; then comes a sketch or a Feynman diagram of
what the section is about; then the equation; then the numbers that the
equation gives for \emph{our} event, every one of them copied from the log
that the standalone script writes in its learning mode
(\cref{sec:ancestral_home:learning_mode}); and finally a plot, drawn from the
same code, that shows the equation at work over the whole phase space rather
than at one point. Equation, number, picture: if one of the three does not
make sense to you, the other two are there to help.

\begin{table}[htbp]\centering\small
\caption[The $q\bar q\to\bbbar$ cross section obtained four ways]{The
$q\bar q\to\bbbar$ cross section at $\sqrt{s}=13.6\TeV$ with
$\pthat\ge100\GeV$ obtained four ways. The first three rows use \PYTHIA's
conventions exactly ($\mb=4.8\GeV$, five incoming flavours, NNPDF3.1 NNLO,
$\muR^2=\muF^2=\pthat^2+\mb^2$, \PYTHIA's own $\alphas$ running); the \MG row
uses that program's defaults for the $b$ mass, the scale and the coupling,
and four incoming flavours (\cref{sec:ancestral_home:madgraph}). The scripts
are in \file{example/ME/Standalone/} and the generator inputs in
\file{example/ME/CMSSW/} and \file{example/ME/MadGraph/} of the JitJet
repository. Quoted errors are statistical or numerical only; the theoretical
uncertainty is the subject of \cref{sec:ancestral_home:scales}.}
\label{tab:ancestral_home:three_numbers}
\begin{tabular}{llr}\toprule
Method & Details & $\sigma$ [pb] \\ \midrule
Standalone script, one pass & LO integral, $2^{16}$ Sobol points & $345.320$ \\
Standalone script, 8 repeats & LO integral, scrambled seeds & $345.322 \pm 0.001$ \\
Full script (LHAPDF, \PYTHIA $\alphas$) & LO integral, 8 repeats & $345.324 \pm 0.001$ \\
\PYTHIA~8.311 in \code{CMSSW\_15\_0\_5} & Monte Carlo, 1000 events & $345.1 \pm 6.1$ \\
\MG~3.8.0 & Monte Carlo, 10\,000 events & $344.3 \pm 1.0$ \\
\bottomrule
\end{tabular}
\end{table}

\paragraph{Why $100\GeV$, and why this process.} A word about the route I
have chosen for us. The first draft of this journey used the more common
generator cut $\pthat\ge30\GeV$. With that cut the two $b$ quarks of our
event both flew off towards the endcaps, one of them into the forward
calorimeter where there is no tracker at all, and a passenger who wants to see
a $b$ jet being tagged would have waited in vain. Raising the cut to
$100\GeV$ costs a factor of fifty in rate, but it buys a mother born at
$|\eta|<1.5$, inside the barrel, where every subdetector of CMS will get a
look at her. Since we follow one event rather than measure a cross section,
that is the better bargain. And a word about the tools: the preface promises
that \MG generates Event~A. For this first stage we use \PYTHIA's own
$2\to2$ matrix element, because it is the process CMS actually runs in the
configuration of \cref{sec:ancestral_home:cmssw}, and because its
leading-order matrix element is a formula short enough to check by hand. \MG
at leading order uses the very same Born matrix element, as
\cref{sec:ancestral_home:madgraph} shows.

\section{Watch the integral happen: the animated journey}
\label{sec:ancestral_home:animation}

We begin, as at every stop of this book, by looking out of the window. The
script \file{example/ME/Standalone/animate_standalone_qqbar_bbbar_xsec.py}
imports the standalone calculation of \cref{sec:ancestral_home:standalone_code}
and renders a four-scene film of about thirty seconds from its tables, its
interpolators and its \code{born\_weights} function, so every number on the
screen comes from the same code that prints the cross section. The film is
on the JitJet playlist,
\begin{center}
  \url{https://www.youtube.com/playlist?list=PLFtvDH9g5Km0}
\end{center}
and \cref{fig:ancestral_home:stills} shows one key frame per scene. Here is
what you will see. Each scene names the equation that explains it; the words
will be new, and that is intended. Watch first, and let the formulae of
\crefrange{sec:ancestral_home:two_protons_meet}{sec:ancestral_home:from_formula_to_number}
catch up with the pictures.
\begin{enumerate}
  \item \textbf{The collision} (\cref{eq:ancestral_home:kinematics}). Our
    event in three dimensions. In the partonic centre-of-mass frame the $u$
    and $\bar u$ collide head on and the $b$ and $\bar b$ leave back to back
    at the angle $\theta^{*}$ to the beam. The translucent tube around the
    beam is the generator cut $\pthat\ge100\GeV$: the quarks must leave
    through its wall, and the sidebar lists $x_1$, $x_2$, $\shat$, $p^{*}$,
    $\cos\theta^{*}$ and $\pthat$ as computed in
    \cref{sec:ancestral_home:matrix_element}. The system is then boosted
    along the beam by the pair rapidity $Y=-2.59$ into the laboratory: the
    tips slide along the tube, because \pt is boost invariant, and the pair
    ends up backward exactly as in \file{event_ME.lhe}.
  \item \textbf{The cube} (\cref{eq:ancestral_home:maps}). Sobol points in
    the unit cube $(u_1,u_2,u_3)$, coloured by $\log_{10}$ of their weight
    $\sum_q w_q$, and then morphed into the physical variables
    $(\log_{10}x_1,\log_{10}x_2,|\cos\theta|)$. Watch the cube become a
    wedge: only the region above the threshold line $x_1x_2=\tau_{\min}$ is
    allowed, and the stretching you see is the Jacobian of
    \cref{sec:ancestral_home:better_variables}, carried in every weight.
  \item \textbf{The surface} (\cref{eq:ancestral_home:weight}). The
    integrand with $\cos\theta$ integrated out,
    $\mathrm{d}\sigma/(\mathrm{d}\log_{10}x_1\,\mathrm{d}\log_{10}x_2)$, as a
    surface over the triangle. Its volume is the cross section. The dome
    sits just above the threshold at $Y=0$, because the PDFs grow towards
    small $x$ and $\hat\sigma\propto1/\shat$, so most of the rate is just
    above the cut; our event, marked on the surface, sits on the backward
    slope at $Y=-2.6$.
  \item \textbf{Convergence} (\cref{eq:ancestral_home:average,eq:ancestral_home:error}).
    One Sobol pass: the cube fills up point by point while the running
    average of the weights, per incoming flavour and in total, settles onto
    $345\pb$, with \PYTHIA's own $345.1\pm6.1\pb$ drawn as a band. The first
    few hundred points swing wildly; by $2^{16}$ the estimate has stopped
    moving at the fourth digit.
\end{enumerate}
To render it yourself you need matplotlib in addition to NumPy and SciPy, and
\code{ffmpeg} for an \code{.mp4} (a \code{.gif} needs only Pillow):
\begin{minted}[linenos=false,frame=single,framesep=1.5mm,fontsize=\small]{bash}
cd example/ME/Standalone
python3 animate_standalone_qqbar_bbbar_xsec.py            # the film, a few minutes
python3 animate_standalone_qqbar_bbbar_xsec.py --fast     # low-resolution preview
python3 animate_standalone_qqbar_bbbar_xsec.py --scene surface -o surface.gif
python3 animate_standalone_qqbar_bbbar_xsec.py --stills stills.pdf --no-video
\end{minted}
The physics options \code{--ecm}, \code{--mb}, \code{--ptmin},
\code{--ptmax}, \code{--power} and \code{--seed} are those of the standalone
script we run next, and \code{--event} points scene~1 at another hard-process
record, so
you can film your own event.

\begin{figure}[htbp]\centering
\includegraphics[page=1,width=0.49\textwidth]{ME/Standalone/animate_standalone_qqbar_bbbar_xsec_stills.pdf}\hfill
\includegraphics[page=2,width=0.49\textwidth]{ME/Standalone/animate_standalone_qqbar_bbbar_xsec_stills.pdf}\\[2mm]
\includegraphics[page=3,width=0.49\textwidth]{ME/Standalone/animate_standalone_qqbar_bbbar_xsec_stills.pdf}\hfill
\includegraphics[page=4,width=0.49\textwidth]{ME/Standalone/animate_standalone_qqbar_bbbar_xsec_stills.pdf}
\caption[Key frames of the animated standalone calculation]{Key frames of
the four scenes of the animated standalone calculation, drawn by
\file{animate_standalone_qqbar_bbbar_xsec.py} with \code{--stills}. Top left:
our event in the partonic frame with the $\pthat\ge100\GeV$ tube. Top right:
Sobol points morphed from the unit cube into
$(\log_{10}x_1,\log_{10}x_2,|\cos\theta|)$. Bottom left: the integrand over
the $(x_1,x_2)$ triangle, whose volume is the cross section. Bottom right:
the running average converging to $345\pb$. The film itself is on the
JitJet YouTube playlist.}
\label{fig:ancestral_home:stills}
\end{figure}

\section{Press the button: the standalone calculation}
\label{sec:ancestral_home:standalone_code}

You have watched the number converge on the screen; now make it converge on
yours. Everything in this book is public. The scripts, configurations and
event records live in the JitJet repository,
\begin{center}
  \url{https://github.com/ravindkv/JitJet}
\end{center}
under \file{example/}, and this chapter's material is in \file{example/ME/}.
I will not print the code here; it would double the length of the chapter,
and a listing on paper cannot be run. Instead, each script gets a short
description, the command to run it, and the output you should see. Run it
now, and read the source on GitHub later, once you have the equations of
this chapter beside you; the function names match the step numbers of
\cref{sec:ancestral_home:from_formula_to_number}, where the integral is
derived.

\subsection{The standalone script}

\file{example/ME/Standalone/standalone_qqbar_bbbar_xsec.py} contains the
calculation of \cref{sec:ancestral_home:from_formula_to_number} and nothing
else. It needs Python~3.10 or newer with NumPy and SciPy; no LHAPDF, no
\PYTHIA. The PDF grid (member~0 of the NNPDF3.1 set on LHAPDF's own knots,
restricted to the phase space of the default run) and \PYTHIA's $\alphas$
curve are embedded as numeric tables between the markers
\code{BEGIN/END GENERATED TABLES}, produced by
\file{dump_standalone_tables.py} in an environment that has both libraries.
A map of the file: \code{TablePDF} (bicubic spline in $(\ln x,\ln Q^2)$,
with the freezing, clipping and heavy-sea conventions of
\cref{sec:ancestral_home:pdfs}); \code{TableAlphaS} (monotone cubic in
$\ln Q^2$, refusing to extrapolate); \code{dsigma\_dcostheta}
(\cref{eq:ancestral_home:dsigma}, vectorised); \code{born\_weights} (steps~1
to~4, that is \cref{eq:ancestral_home:maps,eq:ancestral_home:weight});
\code{integrate} (steps~5 and~6, \cref{eq:ancestral_home:average,eq:ancestral_home:error});
and \code{main} (command line, sanity checks and the printout). From the
repository root:
\begin{minted}[linenos=false,frame=single,framesep=1.5mm,fontsize=\small]{bash}
python3 example/ME/Standalone/standalone_qqbar_bbbar_xsec.py              # one pass
python3 example/ME/Standalone/standalone_qqbar_bbbar_xsec.py --repeats 8  # with error
python3 example/ME/Standalone/standalone_qqbar_bbbar_xsec.py --ptmin 150 --ptmax 300
\end{minted}
The first command takes about a second and ends with
\begin{minted}[linenos=false,frame=single,framesep=1.5mm,fontsize=\small]{text}
TOTAL = 0.345320 +/- nan nb
      = 345.320 +/- nan pb
  d dbar: 0.118568 +/- nan nb
  u ubar: 0.178315 +/- nan nb
  ...
\end{minted}
(the \code{nan} says that one pass has no error estimate). The second gives
$345.322\pm0.001\pb$ after eight passes; the third shows how to select a
window in $\pthat$, which is how sliced production samples are stitched
together. The options are \code{--ecm}, \code{--mb}, \code{--ptmin},
\code{--ptmax}, \code{--power} (the exponent of the number of Sobol points),
\code{--repeats} and \code{--seed}. The two integrators agree to
$5\times10^{-6}$ for the same Sobol seed, so the tables lose nothing; if the
PDF set, the coupling or the kinematic range is changed the tables must be
regenerated, and the script says so.

\subsection{The learning mode: where every number of this chapter comes from}
\label{sec:ancestral_home:learning_mode}

The script has a second voice. With \code{-v} it narrates the calculation
before it performs it: the constants of the integral, then our own event
(read from \file{example/ME/event_ME.lhe}, the record of
\cref{sec:ancestral_home:our_jet}) pushed through \code{born\_weights} one
factor at a time, and finally the first Sobol pass with its running average
after $4, 8, 16, \dots, 65\,536$ points:
\begin{minted}[linenos=false,frame=single,framesep=1.5mm,fontsize=\small]{bash}
python3 example/ME/Standalone/standalone_qqbar_bbbar_xsec.py -v
python3 example/ME/Standalone/standalone_qqbar_bbbar_xsec.py -v --event my_record.lhe
\end{minted}
The output of the first command is saved as
\file{example/ME/Standalone/standalone_qqbar_bbbar_xsec.log}, about a hundred
lines in three numbered sections, and every number quoted in the theory
sections below is taken from it, so that you can put your finger on each
one. Here is the heart of section~[2], the weight of our event's point:
\begin{minted}[linenos=false,frame=single,framesep=1.5mm,fontsize=\footnotesize]{text}
  step 4  the weight (eq. weight)
    Jacobian: ln(1/tau_min) x 2 Y_max x 2 c_max = 8.4367 x 7.0723 x 1.7262 = 102.995
      w_d = 102.995 x 18.916 pb x (1.21237 x 0.00295 + 1.14825 x 0.06629) = 155.268 pb
      w_u = 102.995 x 18.916 pb x (1.25664 x 0.00122 + 1.14546 x 0.21068) = 473.161 pb
      w_s = 102.995 x 18.916 pb x (1.06545 x 0.00060 + 1.06637 x 0.00532) =  12.298 pb
      w_c = 102.995 x 18.916 pb x (0.84207 x 0.00278 + 0.84207 x 0.00278) =   9.114 pb
      w_b = 102.995 x 18.916 pb x (0.60747 x 0.00038 + 0.60747 x 0.00038) =   0.902 pb
    sum_q w_q = 650.742 pb at this point
\end{minted}
You cannot read this yet; by the end of \cref{sec:ancestral_home:from_formula_to_number}
you will be able to reproduce every factor of it by hand. The plots of the
theory sections come from a sibling script in the same directory,
\begin{minted}[linenos=false,frame=single,framesep=1.5mm,fontsize=\small]{bash}
python3 example/ME/Standalone/plot_theory_ME.py        # -> theory_qqbar_bbbar.pdf, one page per section
\end{minted}
which imports the standalone script and uses nothing but its tables and its
\code{born\_weights}; \code{make figures} in the book directory regenerates
the PDF. The only page it cannot draw from the embedded tables is the scale
variation, for which the full script of the next subsection is run through
\file{scan_scales.py}.

\subsection{The same calculation with LHAPDF and \PYTHIA}

The sibling script \file{qqbar_bbbar_xsec.py} in the same directory does the
same integral with LHAPDF for the PDFs and \PYTHIA's own \code{AlphaStrong}
class for the coupling, and offers every knob discussed in this chapter as a
command-line option: upper and lower $\pthat$ and $\hat m$ cuts, $|y|$ or
$|\eta|$ cuts on both quarks, the number of incoming flavours
(\code{--nquark}), all five \PYTHIA scale codes and their multipliers, the
convenience factors \code{--mur-factor} and \code{--muf-factor}, a $K$ factor,
the PDF member, $\alphas(m_Z)$, the coupling backend and order, and the
option to keep negative PDF values (\code{--signed-pdfs}). It needs the
LHAPDF Python bindings and either the \code{pythia8mc} package or an exported
coupling table; its \file{README.md} explains how to build such an environment
and how to export the actual settings of a running \PYTHIA instance with the
header \file{export_pythia_xsec.h}, so that one integrates precisely what the
generator was configured to do rather than what a log file suggests. With
the environment in place:
\begin{minted}[linenos=false,frame=single,framesep=1.5mm,fontsize=\small]{bash}
cd example/ME/Standalone
.venv/bin/python qqbar_bbbar_xsec.py --output result.json
.venv/bin/python qqbar_bbbar_xsec.py --mur-factor 2 --muf-factor 2
.venv/bin/python test_integrator.py
\end{minted}
The first prints $345.324\pm0.001\pb$ and writes the flavour breakdown of
\cref{tab:ancestral_home:flavours}; the second is one row of
\cref{tab:ancestral_home:variations}; the third runs the normalisation and
phase-space checks. We use this script in \cref{sec:ancestral_home:scales}
for the variations that the standalone one does not expose.

\section{Two protons meet: the factorisation theorem in one page}
\label{sec:ancestral_home:two_protons_meet}

You have watched the number appear, and you have made it appear yourself.
Now we earn it. The next five sections derive, factor by factor, the
integral that the script evaluates, end with the seven steps that its
\code{born\_weights} and \code{integrate} functions execute, and close with
an honest estimate of how much the result can be trusted.

Before any formula, a question to guess at. The two protons bring
$13\,600\GeV$ into the collision. How much of it did the hard process of our
event get to use? If your instinct says ``most of it'', look at the record in
\cref{sec:ancestral_home:our_jet}: the pair of $b$ quarks has a mass of
$396\GeV$, under three per cent of the beam energy, and one of the two
colliding partons carried only $15\GeV$. That is the normal case, not an
accident, and the reason is the subject of this section: a proton is not a
particle you can scatter, it is a bag of partons, each carrying a small and
random fraction of the momentum, and the formula below says how to add up
all the ways two of them can meet (\cref{fig:ancestral_home:factorisation_sketch}).

\begin{figure}[htbp]\centering
\resizebox{\textwidth}{!}{% ---------------------------------------------------------------------------
% The factorisation theorem as a picture: two protons, two parton densities,
% one short-distance blob. The x values, energies and the pair mass are those
% of our event (example/ME/event_ME.lhe) as listed in section [2] of the -v log
% of standalone_qqbar_bbbar_xsec.py. Input from chapter_01_ancestral_home.tex
% inside a figure environment.
% ---------------------------------------------------------------------------
\begin{tikzpicture}[x=1cm,y=1cm,
  every node/.style={font=\footnotesize},
  proton/.style={draw,ellipse,minimum width=22mm,minimum height=11mm,thick,fill=gray!10},
  valence/.style={thick,gray!70},
  fermion/.style={thick,postaction={decorate},decoration={markings,mark=at position 0.55 with {\arrow[scale=1.1]{Stealth}}}},
  antifermion/.style={thick,postaction={decorate},decoration={markings,mark=at position 0.45 with {\arrowreversed[scale=1.1]{Stealth}}}},
  remnant/.style={thick,gray!60,postaction={decorate},decoration={markings,mark=at position 0.6 with {\arrow[scale=0.9]{Stealth}}}},
  pdfblob/.style={draw,circle,minimum size=9mm,thick,fill=blue!12},
  hard/.style={draw,circle,minimum size=14mm,very thick,fill=orange!25},
  lab/.style={font=\scriptsize,align=center},
  num/.style={font=\scriptsize,align=center,text=blue!60!black},
  brace/.style={decorate,decoration={brace,amplitude=5pt},thick}]
% the two protons
\node[proton] (PA) at (-5.2,0) {};
\node[lab,above=1pt of PA] {proton A, $+z$, $P=6800\GeV$};
\node[proton] (PB) at (5.2,0) {};
\node[lab,above=1pt of PB] {proton B, $-z$, $P=6800\GeV$};
\foreach \dy in {-0.25,0,0.25} {\draw[valence] ($(PA.west)+(0.15,\dy)$) -- ($(PA.east)+(-0.15,\dy)$);
                                 \draw[valence] ($(PB.west)+(0.15,\dy)$) -- ($(PB.east)+(-0.15,\dy)$);}
% the PDF blobs
\node[pdfblob] (FA) at (-2.6,0) {$f_{\bar u}$};
\node[pdfblob] (FB) at (2.6,0) {$f_{u}$};
\draw[thick,gray!70] (PA.east) -- (FA.west);
\draw[thick,gray!70] (PB.west) -- (FB.east);
\node[lab,below=2pt of FA] {resolved at $\muF$};
\node[lab,below=2pt of FB] {resolved at $\muF$};
% remnants
\draw[remnant] (PA.east) -- (-3.4,-1.5) node[lab,below left=-2pt,text=gray!80!black] {remnant $(1-x_1)P$};
\draw[remnant] (PB.west) -- (3.4,-1.5) node[lab,below right=-2pt,text=gray!80!black] {remnant $(1-x_2)P$};
% the hard blob
\node[hard] (H) at (0,0) {$\hat\sigma$};
\draw[antifermion] (FA.east) -- (H.west);
\draw[fermion] (FB.west) -- (H.east);
\node[num,anchor=south] at (-1.7,0.42) {$\bar u$, $x_1P=14.8\GeV$\\$x_1=2.17\times10^{-3}$};
\node[num,anchor=south] at (1.85,0.42) {$u$, $x_2P=2653\GeV$\\$x_2=0.390$};
% outgoing b and bbar
\draw[fermion] (H.north east) -- (1.5,2.3) node[lab,right,text=red!80!black] {$b$: the mother, $\pt=113\GeV$};
\draw[antifermion] (H.south east) -- (1.5,-2.3) node[lab,right,text=blue!80!black] {$\bar b$, $\pt=113\GeV$};
\node[num,anchor=north east] at (-0.1,-0.9) {$\shat=x_1x_2s$\\$\sqrt{\shat}=396\GeV$};
% what is measured and what is calculated
\draw[brace] (-6.4,1.1) -- (-1.9,1.1) node[midway,above=6pt,lab] {long distance: \emph{measured once},\\the PDF $f_i(x_1,\muF^2)$};
\draw[brace] (1.9,1.1) -- (6.4,1.1) node[midway,above=6pt,lab] {long distance: \emph{measured once},\\the PDF $f_j(x_2,\muF^2)$};
\draw[brace] (0.8,2.75) -- (-0.8,2.75) node[midway,above=6pt,lab] {short distance: \emph{calculated},\\$\hat\sigma_{ij}(\shat;\muR^2,\muF^2)$};
% the sum over flavours and orientations
\node[lab,text=black!65,text width=12.5cm] at (0,-3.3) {$\sigma=\sum_{i,j}\int\!\mathrm{d}x_1\mathrm{d}x_2\,f_i(x_1)f_j(x_2)\,\hat\sigma_{ij}$:
  the sum runs over $d,u,s,c,b$ and both orientations ($q$ from A and $\bar q$ from B, and the reverse);
  our event is the $\bar u$-from-A, $u$-from-B term of the $u$ flavour.};
\end{tikzpicture}
}
\caption[The factorisation theorem as a picture]{The factorisation theorem as a picture. Two protons approach along the beam; from each a single parton is resolved at the scale $\muF$ with a momentum fraction $x$ given by the parton distribution function, and the two partons meet in the short-distance process $\hat\sigma$ that we calculate. The numbers are those of our event, from section~[2] of the \code{-v} log: $x_1=2.17\times10^{-3}$, $x_2=0.390$, $\sqrt{\shat}=396\GeV$. The remnants carry the rest of the two protons away and are forgotten until \cref{ch:whose_child}.}
\label{fig:ancestral_home:factorisation_sketch}
\end{figure}

Look out of the left window: two protons are approaching each other at
$6800\GeV$ apiece. You cannot scatter a proton in perturbation theory; it is
a bound state, and the strong coupling inside it is not small. What rescues
us is the collinear factorisation theorem, which separates the problem into a
long-distance part that is measured once and for all, the parton
distribution functions (PDFs), and a short-distance part that we calculate,
the partonic cross section. For any inclusive hard process
$pp\to\bbbar+X$,
\begin{equation}
\begin{split}
  \sigma(pp\to\bbbar X) &= \sum_{i,j}\int_0^1\!\mathrm{d}x_1\int_0^1\!\mathrm{d}x_2\;
  f_i(x_1,\muF^2)\, f_j(x_2,\muF^2)\;
  \hat\sigma_{ij\to\bbbar}\!\left(\shat;\, \muR^2,\muF^2\right) \\
  &\quad + \mathcal{O}\!\left(\frac{\Lambda^2_{\mathrm{QCD}}}{\pthat^2}\right),
  \qquad \shat = x_1x_2 s .
\end{split}
  \label{eq:ancestral_home:factorisation}
\end{equation}
Read it slowly; it is the most important equation of the chapter. Here
$f_i(x,\muF^2)\,\mathrm{d}x$ is the number of partons of flavour $i$ carrying
a fraction between $x$ and $x+\mathrm{d}x$ of the proton momentum, resolved at
the factorisation scale $\muF$; $\hat\sigma_{ij}$ is the cross section for
the two partons, computed in perturbation theory with the strong coupling
evaluated at the renormalisation scale $\muR$; and $\shat$ is the squared
centre-of-mass energy of the two partons, the only energy the hard scattering
knows about. The two protons are identical, so the sum over $(i,j)$ counts
both a $q$ from beam~1 meeting a $\bar q$ from beam~2 and the reverse. Forget
one orientation and you lose half the cross section; it is the first thing to
check in any integrator.

\paragraph{The numbers.} Section~[2] of the log evaluates the formula at
our event. The $\bar u$ from beam~A carries $x_1=2.174\times10^{-3}$ and the
$u$ from beam~B $x_2=0.3901$, so $\shat=x_1x_2s=156\,895\GeV^2$ and
$\sqrt{\shat}=396.1\GeV$, the mass that the record prints as its momentum
sum. The two orientations of the sum are far from equal at this one point:
the $\bar u$-from-A, $u$-from-B term is $470.2\pb$ of the $u$-flavour weight
and the reverse term $3.0\pb$, because a $u$ quark with $x=2\times10^{-3}$
is common and a $\bar u$ with $x=0.39$ is rare. Integrated over the whole
phase space the two orientations contribute exactly equally, by the symmetry
of the two beams, which is why forgetting one costs a factor of two and not
some odd fraction. \Cref{fig:ancestral_home:factorisation_plot} shows the
same integrand over all of phase space: the rate is concentrated just above
the threshold, $63\%$ of it below $\sqrt{\shat}=305\GeV$, and spread over
five units of pair rapidity; our event, at $396\GeV$ and $Y=-2.59$, sits on
the slope of both distributions, and only $9\%$ of all pairs are boosted
even more than ours.

\begin{figure}[htbp]\centering
\includegraphics[page=1,width=\textwidth]{ME/Standalone/theory_qqbar_bbbar.pdf}
\caption[The integrand of the factorisation formula over the phase space]{The integrand of the factorisation formula, from one Sobol pass of $2^{18}$ points of \file{plot_theory_ME.py}, split by incoming flavour. (a) The distribution in the pair mass $\sqrt{\shat}$: the PDFs rise towards small $x$ and $\hat\sigma\propto1/\shat$, so most of the rate sits just above the threshold $\sqrt{\shat_{\min}}=200\GeV$ set by the $\pthat$ cut. (b) The distribution in the pair rapidity $Y$: flat where both $x$ are small, falling where one parton has to be a valence quark at large $x$. Our event is the dashed line in both panels.}
\label{fig:ancestral_home:factorisation_plot}
\end{figure}

\subsection{Better variables, and a Jacobian you should do once by hand}
\label{sec:ancestral_home:better_variables}

The integral over $x_1$ and $x_2$ is awkward: the PDFs rise steeply towards
small $x$, and the physics depends on the two fractions mostly through their
product. For a $2\to2$ process the natural variables are the pair's
squared-mass fraction and rapidity,
\begin{equation}
  \tau \equiv x_1x_2 = \frac{\shat}{s},\qquad
  Y \equiv \tfrac12\ln\frac{x_1}{x_2},\qquad
  x_{1} = \sqrt{\tau}\,e^{+Y},\quad x_{2} = \sqrt{\tau}\,e^{-Y}.
  \label{eq:ancestral_home:tauY}
\end{equation}
Check the inverse yourself: $x_1x_2=\tau$ and
$\tfrac12\ln(x_1/x_2)=\tfrac12\ln e^{2Y}=Y$. $Y$ is the rapidity of the
$\bbbar$ pair in the laboratory, and it will turn out to be the single number
that decides where in the detector our mother is born.

\begin{figure}[htbp]\centering
\resizebox{\textwidth}{!}{% ---------------------------------------------------------------------------
% The change of variables (x1, x2) -> (ln tau, Y): the allowed triangle, the
% lines of constant tau and Y, and the little rectangle whose image has area
% tau d ln tau dY. Numbers (tau_min, our event) from section [1] and [2] of the
% -v log of standalone_qqbar_bbbar_xsec.py.
% ---------------------------------------------------------------------------
\begin{tikzpicture}[x=1cm,y=1cm,
  every node/.style={font=\footnotesize},
  axis/.style={-{Stealth},thick},
  iso/.style={gray!60,thin},
  lab/.style={font=\scriptsize,align=center},
  head/.style={font=\small\bfseries},
  ev/.style={star,star points=5,star point ratio=2.2,fill=red!80!black,draw=white,inner sep=0pt,minimum size=3.2mm},
  bigarrow/.style={-{Stealth[length=3mm,width=3mm]},gray!60,line width=2.4pt}]
% ---------------- (a) the (ln x1, ln x2) plane: a triangle ---------------------------------
% horizontal: -ln x1 (x1 decreases to the right), vertical: -ln x2; origin = the tau = 1 corner
\begin{scope}
\node[head,anchor=west] at (-0.2,4.9) {(a) the $(\ln x_1,\ln x_2)$ plane};
\fill[blue!8] (0,0) -- (4.2,0) -- (0,4.2) -- cycle;
\draw[axis] (-0.3,0) -- (4.9,0) node[below left] {$-\ln x_1$};
\draw[axis] (0,-0.3) -- (0,4.9) node[below left] {$-\ln x_2$};
\node[lab,anchor=north] at (4.2,-0.1) {$x_1=\tau_{\min}$};
\node[lab,anchor=east] at (-0.1,4.2) {$x_2=\tau_{\min}$};
\node[lab,anchor=north east] at (-0.05,-0.05) {$x_1=x_2=1$\\$(\tau=1)$};
\draw[thick] (4.2,0) -- (0,4.2);
\node[lab,rotate=-45] at (2.6,2.0) {$x_1x_2=\tau_{\min}=2.17\times10^{-4}$};
\foreach \t in {1.4,2.8} {\draw[iso] (\t,0) -- (0,\t);}
\node[lab,rotate=-45,text=gray!80!black] at (1.75,1.75) {$\tau=\mathrm{const}$};
\draw[iso,dashed] (0,0) -- (2.1,2.1);
\draw[iso,dashed] (1.4,0) -- (2.8,1.4);
\draw[iso,dashed] (0,1.4) -- (1.4,2.8);
\node[lab,rotate=45,text=gray!80!black] at (1.3,1.55) {$Y=0$};
\node[lab,rotate=45,text=gray!80!black] at (2.15,0.95) {$Y<0$};
\node[lab,rotate=45,text=gray!80!black] at (0.75,2.35) {$Y>0$};
\node[ev] at (3.05,0.47) {};
\node[lab,anchor=east,text=red!80!black] at (2.9,0.5) {our event, $Y=-2.59$};
\node[lab,text=black!65,text width=4.6cm,anchor=north west] at (-0.2,-0.95) {The region above the threshold line is a triangle: $\tau$ is constant along the anti-diagonals, $Y$ along the diagonals ($Y<0$: $x_1<x_2$, our side).};
\end{scope}
\draw[bigarrow] (5.3,2.2) -- (6.1,2.2);
% ---------------- (b) the (ln tau, Y) plane ---------------------------------------------------
\begin{scope}[xshift=7.2cm]
\node[head,anchor=west] at (-0.2,4.9) {(b) the $(\ln\tau,Y)$ plane};
\fill[blue!8] (0,2.2) -- (4.2,4.3) -- (4.2,0.1) -- cycle;
\draw[axis] (-0.3,2.2) -- (4.9,2.2) node[below left] {$\ln\tau$};
\draw[axis] (0,-0.3) -- (0,4.9) node[below left] {$Y$};
\node[lab,anchor=north] at (0,2.1) {$\ln\tau_{\min}$};
\node[lab,anchor=north] at (4.2,2.1) {$0$};
\draw[thick] (0,2.2) -- (4.2,4.3) node[lab,pos=0.55,above,sloped] {$Y_{\max}=-\tfrac12\ln\tau$};
\draw[thick] (0,2.2) -- (4.2,0.1) node[lab,pos=0.88,below,sloped] {$-Y_{\max}$};
\foreach \t in {1.4,2.8} {\draw[iso] (\t,{2.2-(4.2-\t)/2}) -- (\t,{2.2+(4.2-\t)/2});}
\foreach \y in {1.5,2.9} {\draw[iso,dashed] ({4.2-2*abs(\y-2.2)},\y) -- (4.2,\y);}
\draw[iso,dashed] (0,2.2) -- (4.2,2.2);
% the little rectangle
\fill[orange!60] (1.55,1.5) rectangle (1.95,1.9);
\draw[thick,orange!80!black] (1.55,1.5) rectangle (1.95,1.9);
\node[lab,anchor=south west,text=orange!80!black] at (1.9,1.9) {$\mathrm{d}\ln\tau\,\mathrm{d}Y$};
\node[ev] at (0.68,0.95) {};
\draw[red!80!black,thin] (0.68,0.95) -- (1.6,-0.15);
\node[lab,anchor=north,text=red!80!black] at (1.9,-0.15) {our event: $\ln\tau=-7.07$, $Y=-2.59$ ($73\%$ of $Y_{\max}=3.54$)};
\node[lab,text=black!65,text width=4.6cm,anchor=north west] at (-0.2,-1.15) {The unit cube is mapped linearly onto this triangle: $u_1$ along $\ln\tau$ over the length $\ln(1/\tau_{\min})=8.44$, $u_2$ along $Y$ over $2Y_{\max}(\tau)$.};
\end{scope}
\draw[bigarrow] (12.5,2.2) -- (13.3,2.2);
% ---------------- (c) the Jacobian in the linear (x1, x2) plane --------------------------------
\begin{scope}[xshift=14.4cm]
\node[head,anchor=west] at (-0.2,4.9) {(c) the Jacobian, $\mathrm{d}x_1\mathrm{d}x_2=\tau\,\mathrm{d}\ln\tau\,\mathrm{d}Y$};
\draw[axis] (-0.3,0) -- (4.9,0) node[below left] {$x_1$};
\draw[axis] (0,-0.3) -- (0,4.9) node[below left] {$x_2$};
% a hyperbola tau = const and its neighbour
\draw[iso] plot[domain=0.5:4.3,samples=60] (\x,{2.0/\x});
\draw[iso] plot[domain=0.62:4.3,samples=60] (\x,{2.6/\x});
\node[lab,text=gray!80!black,anchor=west] at (3.4,0.95) {$\tau$ and $\tau e^{\mathrm{d}\ln\tau}$};
% rays Y = const and neighbour
\draw[iso,dashed] (0,0) -- (3.9,3.9);
\draw[iso,dashed] (0,0) -- (4.3,2.9);
\node[lab,text=gray!80!black,anchor=west,rotate=34] at (3.2,2.15) {$Y$ and $Y+\mathrm{d}Y$};
% the parallelogram: corners at the intersections
\coordinate (p1) at (1.414,1.414); \coordinate (p2) at (1.612,1.612);
\coordinate (p3) at (1.964,1.324); \coordinate (p4) at (1.723,1.161);
\fill[orange!60] (p1) -- (p2) -- (p3) -- (p4) -- cycle;
\draw[thick,orange!80!black] (p1) -- (p2) -- (p3) -- (p4) -- cycle;
\node[lab,anchor=south west,text=orange!80!black] at (2.0,1.55) {area $=\tau\,\mathrm{d}\ln\tau\,\mathrm{d}Y$};
\node[lab,text=black!65,text width=4.6cm,anchor=north west] at (-0.2,-0.95) {The rectangle becomes a tilted parallelogram $\tau$ times larger. Since the libraries return $xf$, the two factors of $\tau$ cancel: at our event $f_{\bar u}f_u\,\tau=0.2413=[x_1f_{\bar u}][x_2f_u]$.};
\end{scope}
\end{tikzpicture}
}
\caption[The change of variables from $(x_1,x_2)$ to $(\ln\tau,Y)$]{The change of variables from $(x_1,x_2)$ to $(\ln\tau,Y)$. (a) In the $(\ln x_1,\ln x_2)$ plane the allowed region is a triangle bounded by $x_1x_2=\tau_{\min}$; $\tau$ is constant along one diagonal direction and $Y$ along the other. (b) In the $(\ln\tau,Y)$ plane the same region is the triangle $|Y|\le-\tfrac12\ln\tau$, onto which the unit square of $(u_1,u_2)$ is mapped linearly. (c) A small rectangle $\mathrm{d}\ln\tau\,\mathrm{d}Y$ becomes a parallelogram of area $\tau\,\mathrm{d}\ln\tau\,\mathrm{d}Y$ in the linear $(x_1,x_2)$ plane: the Jacobian of \cref{eq:ancestral_home:jacobian}. Our event is marked with the numbers of section~[2] of the log.}
\label{fig:ancestral_home:jacobian_sketch}
\end{figure}

Now the part I would like you to do once with pencil and paper, because
every Monte Carlo integrator you will ever meet hides it. When we change
variables from $(x_1,x_2)$ to $(\ln\tau,Y)$ the little area element
$\mathrm{d}x_1\,\mathrm{d}x_2$ is stretched, and the Jacobian says by how
much. Write $x_1=e^{\frac12\ln\tau+Y}$ and $x_2=e^{\frac12\ln\tau-Y}$ and
differentiate:
\begin{equation}
  \frac{\partial x_1}{\partial\ln\tau}=\tfrac12x_1,\quad
  \frac{\partial x_1}{\partial Y}=x_1,\qquad
  \frac{\partial x_2}{\partial\ln\tau}=\tfrac12x_2,\quad
  \frac{\partial x_2}{\partial Y}=-x_2 .
\end{equation}
The Jacobian matrix and its determinant are
\begin{equation}
  J=\begin{pmatrix}\tfrac12x_1 & x_1\\[2pt] \tfrac12x_2 & -x_2\end{pmatrix},
  \qquad
  \det J=\left(\tfrac12x_1\right)(-x_2)-x_1\left(\tfrac12x_2\right)=-x_1x_2=-\tau,
\end{equation}
so that
\begin{equation}
  \mathrm{d}x_1\,\mathrm{d}x_2=\tau\,\mathrm{d}\ln\tau\,\mathrm{d}Y .
  \label{eq:ancestral_home:jacobian}
\end{equation}
Picture a tiny rectangle in the $(\ln\tau,Y)$ plane; in the $(x_1,x_2)$
plane it becomes a tilted parallelogram of area $\tau$ times the original.
Nothing deep, but now comes the payoff. PDF libraries do not return $f(x)$;
they return $x f(x,Q^2)$, because that is the quantity of order one across
the whole $x$ range. Multiply and divide by $x_1x_2=\tau$:
\begin{equation}
  f_i(x_1)f_j(x_2)\,\mathrm{d}x_1\mathrm{d}x_2
  = \frac{[x_1f_i(x_1)]\,[x_2f_j(x_2)]}{\tau}\;\tau\,\mathrm{d}\ln\tau\,\mathrm{d}Y
  = [x_1f_i(x_1)]\,[x_2f_j(x_2)]\;\mathrm{d}\ln\tau\,\mathrm{d}Y .
  \label{eq:ancestral_home:xf}
\end{equation}
The two factors of $\tau$ cancel exactly. That is why every integrator in this
chapter samples $\ln\tau$ and $Y$, multiplies the two $xf$ values it gets
from the library, and never divides by $x$. If you ever see a stray factor
$x_1x_2$ in such a code, this is where it came from. The limits are
$\tau_{\min}\le\tau\le1$, with $\tau_{\min}$ fixed by the kinematic cut
(\cref{sec:ancestral_home:matrix_element}), and
\begin{equation}
  |Y|\le Y_{\max}(\tau)\equiv-\tfrac12\ln\tau ,
  \label{eq:ancestral_home:ymax}
\end{equation}
which is simply the statement that both $x_1=\sqrt\tau e^{Y}$ and
$x_2=\sqrt\tau e^{-Y}$ must stay below one. The allowed region in the
$(\ln x_1,\ln x_2)$ plane is a triangle, and you have seen it drawn in
scene~2 of the film (\cref{sec:ancestral_home:animation}). \Cref{app:kinematics} collects these
variables together with the rest of the kinematic vocabulary of the book.

\paragraph{The numbers.} At our event the log gives $\tau=8.483\times10^{-4}$,
$\ln\tau=-7.072$ and $Y=-2.595$; the rapidity range allowed at this $\tau$ is
$|Y|\le Y_{\max}=3.536$, of which our pair uses $73\%$, and the constants of
section~[1] are $\tau_{\min}=2.168\times10^{-4}$ and $\ln(1/\tau_{\min})=8.437$,
the length of the $\ln\tau$ interval. The cancellation of
\cref{eq:ancestral_home:xf} is checked explicitly: the library values are
$x_1f_{\bar u}=1.1455$ and $x_2f_u=0.2107$, the densities themselves are
$f_{\bar u}(x_1)=526.8$ and $f_u(x_2)=0.540$, and
$f_{\bar u}f_u\,\tau=0.24133=[x_1f_{\bar u}][x_2f_u]$ to five digits. The
reason for sampling $\ln\tau$ rather than $\tau$ is in
\cref{fig:ancestral_home:variables_plot}: with the logarithmic map the first
of $24$ bins in $u_1$ carries $26\%$ of the integral and the rest fall off
smoothly, which an integrator can handle; with a linear map in $\tau$ the
first bin would carry more than $99.99\%$, and $96\%$ of the points would be
spent where there is nothing to integrate.

\begin{figure}[htbp]\centering
\includegraphics[page=2,width=\textwidth]{ME/Standalone/theory_qqbar_bbbar.pdf}
\caption[The allowed triangle and the choice of $\ln\tau$ as integration variable]{The change of variables at work, from \file{plot_theory_ME.py}. (a) The integrand $\mathrm{d}^2\sigma/\mathrm{d}\log_{10}x_1\,\mathrm{d}\log_{10}x_2$ over the allowed triangle, with lines of constant $\tau$ and constant $Y$ and our event marked; the colour scale spans five orders of magnitude. (b) The fraction of the integral in each of $24$ bins of $u_1$ when $\ln\tau$ is sampled, as the script does. (c) The same fraction in bins of a coordinate linear in $\tau$: essentially all of the integral falls in the first bin, and a uniform sample would waste almost every point.}
\label{fig:ancestral_home:variables_plot}
\end{figure}

\section{Parton distribution functions: where the mother comes from}
\label{sec:ancestral_home:pdfs}

A question before the formula. A proton is two $u$ quarks and a $d$ quark;
where, then, did the $\bar u$ antiquark of our event come from, and how
likely was it that a proton contained one with a momentum fraction of two
parts in a thousand? Guess a number, then look at the second row of
\cref{tab:ancestral_home:pdf_values}: $x_1f_{\bar u}=1.15$, an antiquark
density \emph{larger} than the valence density $x_2f_u=0.21$ that the $u$
quark of beam~B came with. At small $x$ a proton is mostly gluons and
quark--antiquark pairs that the gluons have split into, and the three valence
quarks are a minority (\cref{fig:ancestral_home:proton_sketch}). The parton
distribution function is the bookkeeping of this crowd.

\begin{figure}[htbp]\centering
\resizebox{\textwidth}{!}{% ---------------------------------------------------------------------------
% What a PDF describes: the proton seen at a resolution mu_F, with the valence
% quarks, the gluons and the sea; our two incoming partons marked with their
% x and xf values (section [2] of the -v log). Panel (b): raising mu_F resolves
% more of the sea and softens the valence, the mechanism behind the mu_F rows.
% ---------------------------------------------------------------------------
\begin{tikzpicture}[x=1cm,y=1cm,
  every node/.style={font=\footnotesize},
  proton/.style={draw,circle,very thick,fill=gray!8},
  val/.style={circle,fill=blue!70!black,inner sep=0pt,minimum size=3.2mm,text=white,font=\tiny},
  sea/.style={circle,fill=violet!70,inner sep=0pt,minimum size=1.8mm},
  seabar/.style={circle,draw=violet!70,thick,fill=white,inner sep=0pt,minimum size=1.8mm},
  gluon/.style={thick,decorate,decoration={coil,aspect=0,segment length=3pt,amplitude=1.3pt},gray!70},
  lab/.style={font=\scriptsize,align=center},
  head/.style={font=\small\bfseries},
  bigarrow/.style={-{Stealth[length=3mm,width=3mm]},gray!60,line width=2.4pt}]
% ---------------- (a) the proton at the scale mu_F = 113 GeV -------------------------------------
\node[head,anchor=west] at (-2.9,3.4) {(a) a proton resolved at $\muF=113\GeV$};
\node[proton,minimum size=52mm] (P) at (0,0) {};
\node[val] (u1) at (-0.9,0.7) {$u$};
\node[val] (u2) at (1.0,0.4) {$u$};
\node[val] (d1) at (0.1,-1.0) {$d$};
\draw[gluon] (u1) -- (u2);
\draw[gluon] (u2) -- (d1);
\draw[gluon] (d1) -- (u1);
\draw[gluon] (u1) -- (-1.9,1.3);
\draw[gluon] (u2) -- (1.9,-0.4);
\draw[gluon] (d1) -- (-0.5,-2.0);
\foreach \px/\py/\qx/\qy in {-1.7/-0.2/-1.45/-0.45,1.5/1.4/1.75/1.15,-0.3/1.9/0.0/1.75,1.2/-1.6/1.45/-1.4,-1.9/1.3/-2.1/1.0,0.7/-2.1/0.45/-2.25}
  {\node[sea] at (\px,\py) {}; \node[seabar] at (\qx,\qy) {};}
\node[lab,anchor=west] at (3.0,2.0) {\tikz\node[val,minimum size=2.6mm]{};\; valence quark: $x\sim0.1$--$0.5$, the proton's identity};
\node[lab,anchor=west] at (3.0,1.3) {\tikz\draw[gluon] (0,0)--(0.6,0);\; gluon: half the momentum, dominant at small $x$};
\node[lab,anchor=west] at (3.0,0.6) {\tikz{\node[sea]{};\node[seabar] at (0.25,0){};}\; sea pair from $g\to q\bar q$: $x\lesssim10^{-2}$, $xf\sim1$};
% our two partons
\draw[-{Stealth},thick,red!70!black] (3.0,-0.5) -- (1.25,0.3);
\node[lab,anchor=west,text=red!70!black] at (3.0,-0.5) {our $u$ (beam B): valence, $x_2=0.390$,\\$x_2f_u=0.2107$, on the broad bump};
\draw[-{Stealth},thick,violet] (3.0,-1.6) -- (-1.3,-0.45);
\node[lab,anchor=west,text=violet] at (3.0,-1.6) {our $\bar u$ (beam A): sea, $x_1=2.17\times10^{-3}$,\\$x_1f_{\bar u}=1.1455$, where $xf$ is of order one};
\node[lab,text=black!65,text width=6.8cm,anchor=north west] at (3.0,-2.3) {$f_i(x,\muF^2)\,\mathrm{d}x$ counts the partons of flavour $i$ with momentum fraction in $[x,x+\mathrm{d}x]$ when the proton is probed at the scale $\muF$. The picture is a snapshot; the numbers are a fit to data.};
% ---------------- (b) the same proton at two resolutions ------------------------------------------
\begin{scope}[xshift=13.4cm]
\node[head,anchor=west] at (-1.4,3.4) {(b) resolution: raising $\muF$ resolves more sea};
\node[proton,minimum size=22mm] (Q1) at (0,1.2) {};
\node[val,minimum size=2.6mm] at (-0.35,1.45) {}; \node[val,minimum size=2.6mm] at (0.4,1.35) {}; \node[val,minimum size=2.6mm] at (0.0,0.75) {};
\node[sea,minimum size=1.5mm] at (-0.6,0.8) {}; \node[seabar,minimum size=1.5mm] at (-0.45,0.65) {};
\node[lab,anchor=west] at (1.3,1.2) {$\muF=57\GeV$ ($Q^2/4$):\\fewer sea quarks, harder valence\\$x_1f_{\bar u}$ lower by $9\%$, $x_2f_u$ higher by $10\%$};
\node[proton,minimum size=22mm] (Q2) at (0,-1.5) {};
\node[val,minimum size=2.6mm] at (-0.35,-1.25) {}; \node[val,minimum size=2.6mm] at (0.4,-1.35) {}; \node[val,minimum size=2.6mm] at (0.0,-1.95) {};
\foreach \px/\py/\qx/\qy in {-0.6/-1.9/-0.45/-2.05,0.6/-0.9/0.75/-1.05,-0.7/-1.1/-0.85/-1.25,0.3/-2.2/0.45/-2.05}
  {\node[sea,minimum size=1.5mm] at (\px,\py) {}; \node[seabar,minimum size=1.5mm] at (\qx,\qy) {};}
\node[lab,anchor=west] at (1.3,-1.5) {$\muF=227\GeV$ ($4Q^2$):\\more sea quarks, softer valence\\$x_1f_{\bar u}$ higher by $8\%$, $x_2f_u$ lower by $9\%$};
\draw[bigarrow] (0,0.0) -- (0,-0.3);
\node[lab,text=black!65,text width=6.6cm,anchor=north west] at (-1.4,-2.9) {DGLAP evolution: the gluon feeds the sea as $Q^2$ grows and the valence quarks lose momentum to radiation. This is why $\muF$ pulls the two PDF factors of our event in opposite directions (\cref{tab:ancestral_home:variations}).};
\end{scope}
\end{tikzpicture}
}
\caption[What a parton distribution function describes]{What a parton distribution function describes. (a) A proton resolved at the scale $\muF=113\GeV$ of our event: three valence quarks, the gluons that bind them and a sea of quark--antiquark pairs from gluon splitting. The $u$ of beam~B is a valence quark at $x_2=0.39$, the $\bar u$ of beam~A a sea antiquark at $x_1=2.2\times10^{-3}$, where $xf$ is of order one. (b) The same proton resolved at half and at twice the scale: raising $\muF$ resolves more of the sea and softens the valence quarks, the DGLAP evolution of \cref{eq:ancestral_home:dglap}. The percentages are those of section~[2] of the \code{-v} log.}
\label{fig:ancestral_home:proton_sketch}
\end{figure}

The PDFs cannot be computed from first principles; they are fitted to data
and evolved in $Q^2$ with the DGLAP equations. CMS Run~3 simulation uses the
NNPDF3.1 NNLO set with $\alphas(m_Z)=0.118$~\cite{NNPDF:2017mvq}, member~0
(the central replica), accessed through LHAPDF~\cite{Buckley:2014ana} under the
name \code{NNPDF31\_nnlo\_as\_0118} and the numeric ID~303600. The set is
distributed as a table of $x f(x,Q^2)$ on a grid of knots in $x$ and $Q$;
LHAPDF interpolates between the knots with cubic splines in $(\ln x,\ln Q^2)$.

\paragraph{The equations.} Three statements define what the table means.
First the definition itself: $f_i(x,\muF^2)\,\mathrm{d}x$ is the number of
partons of flavour $i$ with momentum fraction in $[x,x+\mathrm{d}x]$ when the
proton is probed at the resolution $\muF$. Second the sum rules that pin the
normalisation,
\begin{equation}
  \int_0^1\!\mathrm{d}x\,\bigl[u(x)-\bar u(x)\bigr]=2,\qquad
  \int_0^1\!\mathrm{d}x\,\bigl[d(x)-\bar d(x)\bigr]=1,\qquad
  \sum_i\int_0^1\!\mathrm{d}x\;x f_i(x,\muF^2)=1 ,
  \label{eq:ancestral_home:sumrules}
\end{equation}
which say that the proton has two valence $u$ quarks and one valence $d$
whatever the scale, and that the partons together carry all of its momentum;
the gluon, absent from our process, takes about half. Third the DGLAP
evolution equation, which is the only part of a PDF that perturbation theory
predicts: the dependence on the scale,
\begin{equation}
  \muF^2\,\frac{\partial f_i(x,\muF^2)}{\partial\muF^2}
  =\frac{\alphas(\muF^2)}{2\pi}\sum_j\int_x^1\frac{\mathrm{d}z}{z}\,
  P_{ij}(z)\,f_j\!\left(\frac{x}{z},\muF^2\right).
  \label{eq:ancestral_home:dglap}
\end{equation}
Read it as a flow: when the resolution is raised, a parton $j$ at $x/z$ is
seen to split into a parton $i$ that keeps the fraction $z$ of its momentum,
with the probability $P_{ij}(z)$. The splitting function $P_{qg}$ feeds the
sea from the gluon, so sea densities at small $x$ \emph{rise} with $\muF$;
the function $P_{qq}$ moves a valence quark to smaller $x$ by radiating a
gluon, so the valence density at large $x$ \emph{falls}. Both signs will
matter for us. The $P_{ij}$ themselves are the splitting functions of the
parton shower, and \cref{ch:mother_gives_birth} derives them; here they
are the fit's evolution kernel and the data fix the rest.

Three conventions of the interface matter for reproducing \PYTHIA's numbers,
and you would not guess any of them from the physics:
\begin{itemize}[nosep]
  \item \textbf{Freezing at the grid edges.} Outside the tabulated range in
    $x$ or $Q^2$ the value at the nearest edge is used rather than an
    extrapolation.
  \item \textbf{Clipping.} Fitted PDFs can go slightly negative at large $x$
    and for heavy flavours; \PYTHIA's proton wrapper sets such values to zero.
    In the default integration $1.4$ million of the individual PDF evaluations
    were clipped, yet the effect on the cross section is $0.08\%$
    (\cref{tab:ancestral_home:variations}), because the clipped values are tiny.
  \item \textbf{Symmetric heavy sea.} The charm and bottom antiquark
    densities are taken equal to the quark densities, $\bar c=c$ and
    $\bar b=b$, as in the wrapper.
\end{itemize}

\paragraph{Our event.} Now look at the record of the event we follow through
this book (\cref{sec:ancestral_home:our_jet}). The mother was produced by an
antiquark from beam~A and a quark from beam~B, and the two partons could
hardly be more different. The $\bar u$ carries $x_1=2.17\times10^{-3}$ of its
proton: it is a sea antiquark, drawn from the rising small-$x$ region where
$xf$ is of order one. The $u$ carries $x_2=0.390$: a valence quark, right on
the broad bump of the valence distribution. \Cref{tab:ancestral_home:pdf_values}
gives the values. The asymmetry is the first fact about our mother's birth.
The pair is boosted along the beam with rapidity
$Y=\tfrac12\ln(x_1/x_2)=-2.59$, so both $b$ quarks emerge backward, and the
mother we follow is born at $\eta=-1.44$, comfortably inside the barrel. In
the $30\GeV$ draft the same asymmetry was even stronger ($x_1=2.8\times10^{-4}$,
$Y=-3.47$) and pushed her to the edge of the tracker; the higher cut forces a
larger $\shat=\tau s$, hence a larger $x_1$, hence a smaller boost. That is
the mechanism behind my choice of route.

\begin{table}[htbp]\centering
\caption[PDF values and coupling at our event]{The PDF values and coupling
at our event, evaluated with LHAPDF and with \PYTHIA's $\alphas$ running at
the momentum fractions of \file{example/ME/event_ME.lhe} and at the scale
$Q^2=\pthat^2+\mb^2=12\,842\GeV^2$. \PYTHIA prints the same quantities in the
``Info Listing'' of the \CMSSW log (\cref{sec:ancestral_home:cmssw}); for the
$30\GeV$ draft they agreed with this evaluation to the four printed digits.}
\label{tab:ancestral_home:pdf_values}
\begin{tabular}{llr}\toprule
Parton & $x$ & $xf(x,Q^2)$ \\ \midrule
$\bar u$ from beam A & $2.174\times10^{-3}$ & $1.1455$ \\
$u$ from beam B & $0.3901$ & $0.2107$ \\ \midrule
$\alphas(Q^2)$, \PYTHIA second order & & $0.11425$ \\
$\alphas(Q^2)$ from the PDF set & & $0.11425$ \\
\bottomrule
\end{tabular}
\end{table}

\paragraph{The numbers.} The scale dependence of \cref{eq:ancestral_home:dglap}
is visible in the log, which evaluates the two densities of our event at a
quarter and at four times the scale $Q^2=12\,842\GeV^2$. At $Q^2/4$ the sea
antiquark density is $x_1f_{\bar u}=1.0450$, $8.8\%$ below its central value,
while the valence density is $x_2f_u=0.2322$, $10.2\%$ above; at $4Q^2$ the
two move the other way, by $+8.0\%$ and $-8.5\%$. Their product, which is
what the cross section sees, moves by only $+0.6\%$ and $-1.1\%$: the two
flows of the DGLAP equation nearly cancel at our event, and that is the
origin of the small $\muF$ rows of \cref{tab:ancestral_home:variations}.
\Cref{fig:ancestral_home:pdf_plot} draws the whole table at the event's
scale and the two densities as functions of $Q^2$.

\begin{figure}[htbp]\centering
\includegraphics[page=3,width=\textwidth]{ME/Standalone/theory_qqbar_bbbar.pdf}
\caption[The parton distribution functions of the embedded table]{The parton distribution functions of the embedded table, drawn by \file{plot_theory_ME.py}. (a) $xf(x,Q^2)$ for the five quark flavours and three light antiquarks at the event's scale $Q^2=12\,842\GeV^2$; the gluon is not part of the table because the process does not use it. The two stars are the partons of our event: the sea $\bar u$ at small $x$ and the valence $u$ on the broad bump. (b) The two densities as functions of $Q^2$, relative to their value at the event's scale: the sea rises, the valence falls, and their product is nearly flat, the DGLAP flows of \cref{eq:ancestral_home:dglap} at work.}
\label{fig:ancestral_home:pdf_plot}
\end{figure}

\section{The matrix element for \texorpdfstring{$\bbbar$}{bb} production}
\label{sec:ancestral_home:matrix_element}

\subsection{The \texorpdfstring{$q\bar q\to\bbbar$}{qq->bb} channel at leading order}

How large is a cross section of $30\pb$? Guess before you compute: a
picobarn is $10^{-40}\,\mathrm{m}^2$, so $30\pb$ is a disc of radius
$10^{-19}\,\mathrm{m}$, a ten-thousandth of the proton's size. That is the
area a $u\bar u$ pair of $396\GeV$ presents for turning into a $\bbbar$ pair,
and it is small for two reasons that the one-line formula below makes
explicit: two powers of a coupling of $0.11$, and a $1/\shat$ that says a
hard process is a short-distance process
(\cref{fig:ancestral_home:feynman_qqbar}).

\begin{figure}[htbp]\centering
\resizebox{\textwidth}{!}{% ---------------------------------------------------------------------------
% (a) The one Feynman diagram of q qbar -> b bbar at leading order, with the
% factors that build the matrix element; (b) the kinematics in the partonic
% centre-of-mass frame and the pTHat cut as a tube around the beam. Numbers
% from section [2] of the -v log of standalone_qqbar_bbbar_xsec.py.
% ---------------------------------------------------------------------------
\begin{tikzpicture}[x=1cm,y=1cm,
  every node/.style={font=\footnotesize},
  fermion/.style={thick,postaction={decorate},decoration={markings,mark=at position 0.55 with {\arrow[scale=1.1]{Stealth}}}},
  antifermion/.style={thick,postaction={decorate},decoration={markings,mark=at position 0.45 with {\arrowreversed[scale=1.1]{Stealth}}}},
  gluon/.style={thick,decorate,decoration={coil,aspect=0,segment length=4.5pt,amplitude=2pt}},
  vertex/.style={fill=black,circle,inner sep=1.6pt},
  lab/.style={font=\scriptsize,align=center},
  fac/.style={font=\scriptsize,align=center,text=orange!80!black},
  head/.style={font=\small\bfseries},
  beam/.style={thick,gray!70},
  tube/.style={dashed,red!70!black,thick},
  mom/.style={-{Stealth[length=2.5mm,width=2mm]},very thick}]
% ---------------- (a) the diagram ------------------------------------------------------------------
\node[head,anchor=west] at (-0.3,3.2) {(a) the only diagram: an $s$-channel gluon};
\coordinate (V1) at (1.6,0);
\coordinate (V2) at (4.2,0);
\draw[fermion] (-0.2,1.6) -- (V1) node[pos=0.0,above left,lab] {$q(p_1)$, colour $i$};
\draw[antifermion] (-0.2,-1.6) -- (V1) node[pos=0.0,below left,lab] {$\bar q(p_2)$, colour $j$};
\draw[gluon] (V1) -- (V2);
\draw[fermion] (V2) -- (6.0,1.6) node[above right,lab] {$b(p_3)$, colour $k$, mass $\mb$};
\draw[antifermion] (V2) -- (6.0,-1.6) node[below right,lab] {$\bar b(p_4)$, colour $l$, mass $\mb$};
\node[vertex] at (V1) {}; \node[vertex] at (V2) {};
\node[lab,below] at (2.9,-0.25) {$g^*$, $q^2=(p_1+p_2)^2=\shat$};
\node[fac,below] at (2.9,-0.7) {propagator $1/\shat$};
\node[fac,anchor=south] at (1.6,0.35) {$g_{\mathrm s}\,T^a_{ji}$};
\node[fac,anchor=south] at (4.2,0.35) {$g_{\mathrm s}\,T^a_{kl}$};
\node[lab,text=black!65,text width=6.6cm,anchor=north west] at (-0.3,-2.1) {$|\mathcal M|^2\propto g_{\mathrm s}^4=16\pi^2\alphas^2$; colour: $\tfrac19\sum|T^a_{ji}T^a_{kl}|^2=\tfrac19\cdot\mathrm{Tr}(T^aT^b)\mathrm{Tr}(T^aT^b)=\tfrac29$; spin average $\tfrac14$; angular structure $1+\cos^2\theta$ from a spin-1 exchange between two spin-$\tfrac12$ pairs. The result is \cref{eq:ancestral_home:dsigma}.};
% ---------------- (b) the kinematics in the partonic frame -----------------------------------------
\begin{scope}[xshift=9.4cm]
\node[head,anchor=west] at (-0.3,3.2) {(b) the partonic frame and the $\pthat\ge100\GeV$ tube};
\draw[beam] (-0.3,0) -- (7.3,0) node[right,lab] {beam ($z$)};
\draw[tube] (-0.3,1.3) -- (7.3,1.3) node[right,lab,text=red!70!black] {$\pthat^{\min}$};
\draw[tube] (-0.3,-1.3) -- (7.3,-1.3);
\coordinate (O) at (3.5,0);
% incoming
\draw[mom,green!50!black] (1.0,0) -- (3.1,0) node[midway,above,lab] {$u$, $\sqrt{\shat}/2$};
\draw[mom,violet] (6.0,0) -- (3.9,0) node[midway,above,lab] {$\bar u$};
% outgoing b (red) and bbar (blue), angle theta* to +z; cos theta*(bbar) = -0.82
\coordinate (B) at ($(O)+(-2.1,1.47)$);
\coordinate (Bb) at ($(O)+(2.1,-1.47)$);
\draw[mom,blue!80!black] (O) -- (B); \node[anchor=south,lab,text=blue!80!black] at (B) {$\bar b$: $p^*=198\GeV$, $\cos\theta^*=-0.820$};
\draw[mom,red!80!black] (O) -- (Bb) node[below right,lab,text=red!80!black] {$b$: $p^*=198\GeV$};
\draw[-{Stealth},thin] (O) ++(0.9,0) arc[start angle=0,end angle=145,radius=0.9];
\node[lab] at ($(O)+(0.25,1.15)$) {$\theta^*=2.53$};
% pT as the projection
\draw[dashed,thin] (B) -- ($(O)+(-2.1,0)$);
\node[lab,anchor=east] at ($(O)+(-2.2,0.7)$) {$\pt=p^*\sin\theta^*$\\$=113\GeV$};
% the allowed cone: |cos theta| <= c_max
\fill[red!15,opacity=0.6] (O) -- ($(O)+(2.6,1.3)$) -- ($(O)+(2.6,0)$) -- cycle;
\fill[red!15,opacity=0.6] (O) -- ($(O)+(2.6,-1.3)$) -- ($(O)+(2.6,0)$) -- cycle;
\fill[red!15,opacity=0.6] (O) -- ($(O)+(-2.6,1.3)$) -- ($(O)+(-2.6,0)$) -- cycle;
\fill[red!15,opacity=0.6] (O) -- ($(O)+(-2.6,-1.3)$) -- ($(O)+(-2.6,0)$) -- cycle;
\node[lab,text=red!70!black] at ($(O)+(2.1,0.45)$) {forbidden:\\$|\cos\theta|>c_{\max}$};
\node[lab,text=black!65,text width=7.4cm,anchor=north west] at (-0.3,-1.7) {Both quarks must leave through the wall of the tube, never through its ends: $|\cos\theta|\le c_{\max}=\sqrt{1-(\pthat^{\min}/p^*)^2}=0.863$ at our $\shat$, which removes $14\%$ of the solid angle and $19\%$ of $\hat\sigma$. The whole frame is then boosted along $z$ by $Y=-2.59$; \pt does not change.};
\end{scope}
\end{tikzpicture}
}
\caption[The leading-order diagram and the kinematics of the partonic frame]{The leading-order process. (a) The only Feynman diagram of $q\bar q\to\bbbar$ at this order, an $s$-channel gluon, with the factors that build \cref{eq:ancestral_home:dsigma}: a coupling $g_{\mathrm s}T^a$ at each vertex, the propagator $1/\shat$, and the colour and spin averages. (b) The kinematics in the partonic centre-of-mass frame: the quarks leave back to back at the angle $\theta^*$ to the beam with momentum $p^*$, and the generator cut $\pthat\ge100\GeV$ is a tube around the beam through whose wall both must leave, $|\cos\theta|\le c_{\max}$ (\cref{eq:ancestral_home:cmax}). Numbers from section~[2] of the \code{-v} log.}
\label{fig:ancestral_home:feynman_qqbar}
\end{figure}

Here is the engine of the whole stop, and it fits on one line. The Born
process is a single $s$-channel gluon exchange. Averaging over the spins and
colours of the incoming massless quark pair and summing over those of the
massive outgoing pair, the differential cross section in the partonic
centre-of-mass frame is~\cite{ParticleDataGroup:2024cfk}
\begin{equation}
  \frac{\mathrm{d}\hat\sigma}{\mathrm{d}\cos\theta}
  = \frac{\pi\alphas^2(\muR^2)\,\beta}{9\,\shat}
    \left[\,1+\cos^2\theta+\rho\,(1-\cos^2\theta)\,\right],
  \qquad
  \rho\equiv\frac{4\mb^2}{\shat},\quad \beta\equiv\sqrt{1-\rho},
  \label{eq:ancestral_home:dsigma}
\end{equation}
where $\theta$ is the angle between the incoming quark and the outgoing $b$,
$\beta$ is the $b$ velocity in that frame and $\rho$ vanishes far above
threshold. Let me point at each factor as we pass it. The $\alphas^2$ is two
strong vertices. The $1/\shat$ is dimensional analysis: a cross section is an
area, and $\shat$ is the only scale left once $\rho\to0$. The $\beta$ is
phase space: two heavy quarks near threshold have nowhere to go. The
$1+\cos^2\theta$ is the angular distribution of any spin-$\tfrac12$ pair
produced through a spin-1 exchange; you have met it in $e^+e^-\to\mu^+\mu^-$,
and QCD only changes the prefactor. And the prefactor $1/9$ is the colour
factor: the average $\tfrac13\cdot\tfrac13$ over the incoming colours, the
sum $\mathrm{Tr}(T^aT^b)\,\mathrm{Tr}(T^aT^b)=\tfrac14\delta^{ab}\delta^{ab}=2$
over the outgoing colours and the exchanged gluon, and the spin average, which
together turn the QED result into this one.

Integrating over $\cos\theta$ from $-1$ to $1$ is a two-line job, and worth
doing because the standalone script checks itself against the result:
\begin{equation}
  \int_{-1}^{1}(1+c^2)\,\mathrm{d}c=\frac83,\qquad
  \int_{-1}^{1}(1-c^2)\,\mathrm{d}c=\frac43,\qquad\Longrightarrow\qquad
  \hat\sigma(\shat) = \frac{\pi\alphas^2\beta}{9\shat}\left[\frac83+\frac{4\rho}{3}\right]
  = \frac{4\pi\alphas^2\,\beta\,(2+\rho)}{27\,\shat},
  \label{eq:ancestral_home:sigmahat}
\end{equation}
which reduces to the familiar $8\pi\alphas^2/(27\shat)$ for massless quarks.
This is exactly the function \code{sigmaKin} of \PYTHIA's class
\code{Sigma2qqbar2QQbar}, process code~124. At our event,
$\shat=156\,895\GeV^2$ and $\alphas=0.1143$ give $\hat\sigma=30.2\pb$ and
$\mathrm{d}\hat\sigma/\mathrm{d}\cos\theta=18.9\pb$ at the actual
$\cos\hat\theta=-0.820$; you are asked to verify both in Exercise~2. The
conversion from natural units is
$(\hbar c)^2 = 0.389\,379\,372\GeV^2\,\mathrm{mb} = 3.893\,793\,72\times10^{8}\GeV^2\,\mathrm{pb}$.
Section~[2] of the log prints both numbers to more digits, from the very
function the integrator calls: $\mathrm{d}\hat\sigma/\mathrm{d}\cos\theta=18.916\pb$
at $\cos\theta^*=-0.8203$, and $\hat\sigma=30.157\pb$, with $\rho=5.87\times10^{-4}$
and $\beta=0.99971$: at this energy the $b$ mass is invisible in the matrix
element, and the massless $8\pi\alphas^2/(27\shat)$ would be off by three
parts in ten thousand.

\paragraph{Kinematics of the two quarks.} In the partonic frame the two
quarks share the energy equally, $E_b=\sqrt{\shat}/2$, so each has momentum
\begin{equation}
  p^{*}=\sqrt{E_b^2-\mb^2}=\sqrt{\frac{\shat}{4}-\mb^2},
  \label{eq:ancestral_home:pstar}
\end{equation}
and the transverse momentum is the component perpendicular to the beam,
$\pt=p^{*}\sin\theta$. Squaring,
\begin{equation}
  \pthat^2 = p^{*2}\,(1-\cos^2\theta),\qquad
  y^{*} = \operatorname{artanh}(\beta\cos\theta),\qquad
  y_{b,\bar b} = Y \pm y^{*},\qquad
  \eta = \operatorname{arsinh}\!\left(\frac{m_{\mathrm{T}}}{\pt}\sinh y\right),
  \label{eq:ancestral_home:kinematics}
\end{equation}
with $m_{\mathrm{T}}^2=\pt^2+\mb^2$. Rapidities add under boosts along the
beam, which is why the lab rapidities of the two quarks are just $Y$ plus and
minus the rapidity $y^{*}$ each has in the partonic frame
(\cref{app:kinematics} proves this). At Born level the two quarks have equal
and opposite transverse momentum, so ``$\pthat$'' is the \pt of either.

Now watch what the generator cut does. A cut $\pthat\ge\pthat^{\min}$
requires $p^{*2}(1-\cos^2\theta)\ge(\pthat^{\min})^2$, that is
\begin{equation}
  |\cos\theta|\le c_{\max}\equiv\sqrt{1-\left(\frac{\pthat^{\min}}{p^{*}}\right)^2},
  \label{eq:ancestral_home:cmax}
\end{equation}
so the cut on \pt is a cut on the angle: the quarks must leave through the
wall of a tube of radius $\pthat^{\min}$ around the beam. It can only be
satisfied when $p^{*}>\pthat^{\min}$, that is when
\begin{equation}
  \shat\ge\shat_{\min}=4\left(\mb^2+(\pthat^{\min})^2\right),\qquad
  \tau_{\min}=\frac{\shat_{\min}}{s}=\frac{4\,(4.8^2+100^2)}{13600^2}=2.17\times10^{-4}.
  \label{eq:ancestral_home:taumin}
\end{equation}
This is the $\tau_{\min}$ of \cref{sec:ancestral_home:better_variables}. At
our event $p^{*}=198\GeV$ and $c_{\max}=0.863$: the cut removes the
$14\%$ of the solid angle nearest the beam, which carries $19\%$ of the
partonic cross section because of the $1+\cos^2\theta$ shape. Note finally
that rapidity $y$ and pseudorapidity $\eta$ differ for a massive quark; the
difference is $0.001$ at $\pt=113\GeV$ and would be far larger near
threshold.
The log confirms the three numbers of the cut at our event, $c_{\max}=0.8631$,
a removed solid-angle fraction of $0.137$ and a surviving fraction of
$\hat\sigma$ of $0.808$, and it recovers the record's own kinematics from
the formulae: $p^*=197.99\GeV$ against $197.995$ from the boosted
four-vectors, $\pthat=113.23\GeV$ against $113.22$, and the two rapidities
$y_b=-1.438$ and $y_{\bar b}=-3.752$. \Cref{fig:ancestral_home:matrix_element_plot}
shows the three functions of this subsection over their whole range.

\begin{figure}[htbp]\centering
\includegraphics[page=4,width=\textwidth]{ME/Standalone/theory_qqbar_bbbar.pdf}
\caption[The Born matrix element and the effect of the $\pthat$ cut]{The Born matrix element and what the generator cut does to it, from \file{plot_theory_ME.py}. (a) The angular distribution \cref{eq:ancestral_home:dsigma} at the $\shat$ of our event, with the regions removed by $|\cos\theta|>c_{\max}$ shaded and our event's $\cos\theta^*$ marked. (b) The integrated $\hat\sigma(\shat)$ of \cref{eq:ancestral_home:sigmahat} at fixed $\alphas$ from the $2\mb$ threshold, where $\beta\to0$, to $3\TeV$; the massless form is indistinguishable above a few tens of GeV. (c) The cut as a cut on the angle: $c_{\max}$ and the fraction of $\hat\sigma$ that survives, as functions of $\sqrt{\shat}$; near the threshold $\sqrt{\shat_{\min}}=200\GeV$ almost everything is removed, at our event $81\%$ survives.}
\label{fig:ancestral_home:matrix_element_plot}
\end{figure}

\subsection{The \texorpdfstring{$gg$}{gg} channel, the incoming \texorpdfstring{$b$}{b}, and four versus five flavours}
\label{sec:ancestral_home:gg_and_flavours}

Which quark flavour makes most of our $\bbbar$ pairs, and which makes the
fewest? Most readers guess $u\bar u$ first, correctly, because the proton has
two valence $u$ quarks; fewer guess that $d\bar d$ is a third of the total
rather than a half of the $u\bar u$ share, and almost nobody expects an
incoming \emph{bottom} quark to contribute at all. It does, at the per cent
level, and whether it should is a question of convention that this
subsection has to settle before the number can be compared with anything
(\cref{fig:ancestral_home:feynman_gg}).

\begin{figure}[htbp]\centering
\resizebox{\textwidth}{!}{% ---------------------------------------------------------------------------
% What is left out and why: (a) the three gg -> b bbar diagrams that
% HardQCD:gg2bbbar = off removes, (b) the two b bbar -> b bbar diagrams of
% which PYTHIA process 124 keeps only the first, (c) the five-flavour picture
% of a b quark inside the proton (g -> b bbar in the evolution) against the
% four-flavour one, where the same splitting sits in the matrix element.
% ---------------------------------------------------------------------------
\begin{tikzpicture}[x=1cm,y=1cm,
  every node/.style={font=\footnotesize},
  fermion/.style={thick,postaction={decorate},decoration={markings,mark=at position 0.55 with {\arrow[scale=1.0]{Stealth}}}},
  antifermion/.style={thick,postaction={decorate},decoration={markings,mark=at position 0.45 with {\arrowreversed[scale=1.0]{Stealth}}}},
  plain/.style={thick},
  gluon/.style={thick,decorate,decoration={coil,aspect=0,segment length=4pt,amplitude=1.8pt}},
  vertex/.style={fill=black,circle,inner sep=1.4pt},
  lab/.style={font=\scriptsize,align=center},
  head/.style={font=\small\bfseries},
  off/.style={draw=red!60!black,dashed,rounded corners=3pt,thick},
  on/.style={draw=green!50!black,rounded corners=3pt,thick}]
% ---------------- (a) gg -> b bbar: switched off ---------------------------------------------------
\node[head,anchor=west] at (-0.3,2.6) {(a) $gg\to\bbbar$: three diagrams, switched off in our run};
\begin{scope}[xshift=0cm]
\draw[gluon] (0,1.2) -- (1.2,0.6); \draw[gluon] (0,-1.2) -- (1.2,-0.6);
\draw[gluon] (1.2,0.6) -- (1.2,-0.6);
\draw[fermion] (1.2,0.6) -- (2.4,1.2) node[right,lab] {$b$}; \draw[antifermion] (1.2,-0.6) -- (2.4,-1.2) node[right,lab] {$\bar b$};
\node[vertex] at (1.2,0.6) {}; \node[vertex] at (1.2,-0.6) {};
\node[lab,below] at (1.2,-1.4) {$t$-channel $b$};
\end{scope}
\begin{scope}[xshift=3.6cm]
\draw[gluon] (0,1.2) -- (1.2,0.6); \draw[gluon] (0,-1.2) -- (1.2,-0.6);
\draw[gluon] (1.2,0.6) -- (1.2,-0.6);
\draw[fermion] (1.2,0.6) -- (2.4,-1.2) node[right,lab] {$b$}; \draw[antifermion] (1.2,-0.6) -- (2.4,1.2) node[right,lab] {$\bar b$};
\node[vertex] at (1.2,0.6) {}; \node[vertex] at (1.2,-0.6) {};
\node[lab,below] at (1.2,-1.4) {$u$-channel $b$};
\end{scope}
\begin{scope}[xshift=7.2cm]
\draw[gluon] (0,1.2) -- (0.8,0); \draw[gluon] (0,-1.2) -- (0.8,0);
\draw[gluon] (0.8,0) -- (1.8,0);
\draw[fermion] (1.8,0) -- (2.8,1.2) node[right,lab] {$b$}; \draw[antifermion] (1.8,0) -- (2.8,-1.2) node[right,lab] {$\bar b$};
\node[vertex] at (0.8,0) {}; \node[vertex] at (1.8,0) {};
\node[lab,below] at (1.4,-1.4) {$s$-channel $g$,\\three-gluon vertex};
\end{scope}
\draw[off] (-0.4,-2.4) rectangle (10.4,1.6);
\node[lab,text=red!60!black,anchor=north] at (5.0,-2.5) {\code{HardQCD:gg2bbbar = off}: dominant in an inclusive $\bbbar$ sample, and not a one-line formula};
% ---------------- (b) b bbar -> b bbar: only the s-channel is used -----------------------------------
\begin{scope}[xshift=11.6cm]
\node[head,anchor=west] at (-0.3,2.6) {(b) the incoming $b$: $b\bar b\to b\bar b$};
\begin{scope}
\draw[fermion] (0,1.2) -- (0.8,0); \draw[antifermion] (0,-1.2) -- (0.8,0);
\draw[gluon] (0.8,0) -- (1.8,0);
\draw[fermion] (1.8,0) -- (2.6,1.2); \draw[antifermion] (1.8,0) -- (2.6,-1.2);
\node[vertex] at (0.8,0) {}; \node[vertex] at (1.8,0) {};
\node[lab,below] at (1.3,-1.4) {$s$-channel: used\\(process 124,\\all five flavours)};
\draw[on] (-0.3,-2.4) rectangle (2.9,1.6);
\end{scope}
\begin{scope}[xshift=4.2cm]
\draw[fermion] (0,1.2) -- (1.0,0.6) -- (2.2,1.2); \draw[antifermion] (0,-1.2) -- (1.0,-0.6) -- (2.2,-1.2);
\draw[gluon] (1.0,0.6) -- (1.0,-0.6);
\node[vertex] at (1.0,0.6) {}; \node[vertex] at (1.0,-0.6) {};
\node[lab,below] at (1.1,-1.4) {$t$-channel: ignored\\(part of\\the $1.6\%$)};
\draw[off] (-0.3,-2.4) rectangle (2.5,1.6);
\end{scope}
\end{scope}
% ---------------- (c) five against four flavours -------------------------------------------------------
\begin{scope}[yshift=-6.2cm]
\node[head,anchor=west] at (-0.3,2.3) {(c) where the $b$ of the proton comes from: five-flavour against four-flavour scheme};
\begin{scope}
\node[draw,ellipse,thick,fill=gray!10,minimum width=16mm,minimum height=8mm] (P) at (0.6,0) {$p$};
\draw[gluon] (P.east) -- (2.4,0);
\draw[fermion] (2.4,0) -- (3.6,0.9) node[right,lab] {$b$ enters $\hat\sigma$ (massless)};
\draw[antifermion] (2.4,0) -- (3.6,-0.9) node[right,lab] {$\bar b$ stays in the remnant};
\node[vertex] at (2.4,0) {};
\node[lab,text=black!65,text width=7.4cm,anchor=north west] at (-0.3,-1.3) {5FS (\PYTHIA, CMS, the scripts): the $g\to\bbbar$ splitting is inside the PDF $f_b(x,\muF^2)$, with its $\ln(\muF^2/\mb^2)$ resummed by DGLAP; $f_b=0$ below $\mb$. Incoming-$b$ share here: $1.6\%$.};
\end{scope}
\begin{scope}[xshift=11.6cm]
\node[draw,ellipse,thick,fill=gray!10,minimum width=16mm,minimum height=8mm] (P) at (0.6,0) {$p$};
\draw[gluon] (P.east) -- (2.4,0);
\draw[fermion] (2.4,0) -- (3.6,0.9) node[right,lab] {$b$ in the matrix element};
\draw[antifermion] (2.4,0) -- (3.6,-0.9) node[right,lab] {$\bar b$ in the final state};
\node[vertex] at (2.4,0) {};
\node[draw,circle,thick,fill=orange!25,minimum size=7mm] at (4.9,0.0) {$\hat\sigma$};
\node[lab,text=black!65,text width=7.4cm,anchor=north west] at (-0.3,-1.3) {4FS (\MG's default proton): no $b$ in the proton; the same splitting is an explicit $2\to3$ (or higher) matrix element with the $b$ mass kept. Equal to 5FS once higher orders are included; $1.6\%$ lower at this order.};
\end{scope}
\end{scope}
\end{tikzpicture}
}
\caption[What is left out: the $gg$ channel, the $t$-channel $b\bar b$ diagram, and the two flavour schemes]{What is left out of the calculation and why. (a) The three diagrams of $gg\to\bbbar$, switched off in our configuration; they dominate an inclusive $\bbbar$ sample. (b) For an incoming $b$ the process $b\bar b\to b\bar b$ has an $s$-channel and a $t$-channel diagram; \PYTHIA process~124, and our integrators, keep only the first. (c) The five-flavour scheme puts the $g\to\bbbar$ splitting into the PDF $f_b$ and resums its logarithm; the four-flavour scheme keeps the $b$ out of the proton and puts the same splitting into the matrix element. The $1.6\%$ is the incoming-$b$ share of section~[3] of the \code{-v} log.}
\label{fig:ancestral_home:feynman_gg}
\end{figure}

At LHC energies the gluon luminosity is so large that $gg\to\bbbar$ dominates
the inclusive $\bbbar$ rate. Our example switches it off
(\code{HardQCD:gg2bbbar = off}) to keep a single channel whose matrix element
fits on one line. Nothing in the rest of the journey depends on this choice,
but a passenger comparing with an inclusive $\bbbar$ sample must remember it.

The sum over $q$ in \cref{eq:ancestral_home:factorisation} runs over the
flavours \PYTHIA allows in the initial state, \code{PDFinProcess:nQuarkIn=5}
by default: $d,u,s,c$ and $b$. For $q=b$ the process $b\bar b\to b\bar b$ has
additional $t$-channel diagrams that \cref{eq:ancestral_home:dsigma} ignores;
\PYTHIA process~124 applies the $s$-channel formula to all five flavours, and
so do our integrators. This is a convention we reproduce, not a complete
calculation. It matters at the $1.6\%$ level, the size of the incoming-$b$
term (\cref{tab:ancestral_home:flavours}).

Written out, the flavour sum of \cref{eq:ancestral_home:factorisation} for
our process is
\begin{equation}
  \sigma=\sum_{q=d,u,s,c,b}\sigma_q,\qquad
  \sigma_q=\int\!\mathrm{d}x_1\mathrm{d}x_2\,
  \bigl[f_q(x_1)f_{\bar q}(x_2)+f_{\bar q}(x_1)f_q(x_2)\bigr]\,
  \hat\sigma_{q\bar q\to\bbbar}(\shat),
  \label{eq:ancestral_home:flavoursum}
\end{equation}
with the \emph{same} $\hat\sigma$ for every flavour; the flavours differ only
through their PDFs. The shares are therefore a statement about the proton,
not about the matrix element. Section~[3] of the log lists them for the
default run: $u\bar u$ $51.6\%$, $d\bar d$ $34.3\%$, $s\bar s$ $8.8\%$,
$c\bar c$ $3.7\%$ and $b\bar b$ $1.6\%$. The $u$-to-$d$ ratio of $1.5$ rather
than $2$ is the valence-quark count diluted by the sea, which is
flavour-symmetric to a first approximation; and the heavy-quark shares are
nothing but the heavy sea of \cref{fig:ancestral_home:pdf_plot}, which is
why the $b\bar b$ term exists at all. At our event's own point the weights
are $w_u=473.2$, $w_d=155.3$, $w_s=12.3$, $w_c=9.1$ and $w_b=0.9\pb$, more
$u$-heavy than the average because the valence $u$ at $x_2=0.39$ is three times as
abundant as the valence $d$ there, by the $xf$ values of the log. \Cref{fig:ancestral_home:flavours_plot}
follows the shares as the generator cut is raised: the valence $u$ wins
steadily because a higher cut forces larger $x$, and the incoming $b$ falls
from two per cent at $30\GeV$ to two per mille at $800\GeV$.

Treating the $b$ as a massless parton inside the proton while giving it a
mass in the final state is the \emph{five-flavour scheme} (5FS). The
alternative \emph{four-flavour scheme} (4FS) keeps the $b$ out of the proton
and generates it only in the matrix element, at the price of large logarithms
$\ln(\pthat^2/\mb^2)$ that the 5FS resums through PDF evolution. For $\bbbar$
production at $\pthat\ge100\GeV$ the two schemes must agree once higher orders
are included; at leading order they do not, and the 5FS with
\cref{eq:ancestral_home:dsigma} is simply what \PYTHIA and CMS use. \MG, as
we will see, takes the 4FS road by default.


\begin{figure}[htbp]\centering
\includegraphics[page=5,width=\textwidth]{ME/Standalone/theory_qqbar_bbbar.pdf}
\caption[The share of each incoming flavour as a function of the generator cut]{The share of each incoming flavour, from single Sobol passes of the standalone script at eleven values of the generator cut (\file{plot_theory_ME.py}). (a) The stacked shares: the valence $u$ grows from $45\%$ at $30\GeV$ to $63\%$ at $800\GeV$, the sea flavours shrink, and the incoming $b$ falls from $2.2\%$ to $0.2\%$. The percentages at $100\GeV$ are those of \cref{tab:ancestral_home:flavours}. (b) The total cross section against the cut, with the local power $n$ of $\sigma\propto(\pthat^{\min})^{-n}$ between neighbouring points: dimensional analysis alone would give $n=4$, the falling PDFs and the running coupling make it steeper at high \pt, and the cut on an angle rather than on an energy makes it flatter at low \pt.}
\label{fig:ancestral_home:flavours_plot}
\end{figure}

\subsection{Scales and the coupling}
\label{sec:ancestral_home:scales_and_coupling}

Here is a question with no right answer, which is exactly the point. The
hard process of our event involves a $b$ mass of $4.8\GeV$, a transverse
momentum of $113\GeV$, a pair mass of $396\GeV$ and a beam energy of
$13\,600\GeV$. At which of these should the strong coupling be evaluated?
QCD does not say; it only says that the full answer, to all orders, does not
depend on the choice, and that a leading-order answer does
(\cref{fig:ancestral_home:scales_sketch}). Guess how much it depends,
then read on.

\begin{figure}[htbp]\centering
\resizebox{\textwidth}{!}{% ---------------------------------------------------------------------------
% Where the two scales enter and how far apart the scales of the problem are.
% (a) the factorised picture with mu_F at the parton densities and mu_R at
% the vertices; (b) the scales of our event on a logarithmic ladder with the
% coupling from the -v log of standalone_qqbar_bbbar_xsec.py (section [2],
% step 3) and the value at mZ from the table.
% ---------------------------------------------------------------------------
\begin{tikzpicture}[x=1cm,y=1cm,
  every node/.style={font=\footnotesize},
  fermion/.style={thick,postaction={decorate},decoration={markings,mark=at position 0.55 with {\arrow[scale=1.1]{Stealth}}}},
  antifermion/.style={thick,postaction={decorate},decoration={markings,mark=at position 0.45 with {\arrowreversed[scale=1.1]{Stealth}}}},
  gluon/.style={thick,decorate,decoration={coil,aspect=0,segment length=4.5pt,amplitude=2pt}},
  pdfblob/.style={draw,circle,minimum size=9mm,thick,fill=blue!12},
  vertex/.style={fill=black,circle,inner sep=1.6pt},
  muF/.style={draw=blue!60!black,dashed,thick,rounded corners=3pt},
  muR/.style={draw=orange!80!black,dashed,thick,rounded corners=3pt},
  lab/.style={font=\scriptsize,align=center},
  head/.style={font=\small\bfseries},
  tick/.style={thick},
  ladder/.style={-{Stealth},very thick}]
% ---------------- (a) where mu_F and mu_R enter ----------------------------------------------------
\node[head,anchor=west] at (-0.3,3.0) {(a) two conventions, two scales};
\node[pdfblob] (FA) at (0.3,1.1) {$f_{\bar u}$};
\node[pdfblob] (FB) at (0.3,-1.1) {$f_{u}$};
\coordinate (V1) at (2.6,0); \coordinate (V2) at (4.6,0);
\draw[antifermion] (FA.east) -- (V1);
\draw[fermion] (FB.east) -- (V1);
\draw[gluon] (V1) -- (V2);
\draw[fermion] (V2) -- (6.4,1.3) node[right,lab] {$b$};
\draw[antifermion] (V2) -- (6.4,-1.3) node[right,lab] {$\bar b$};
\node[vertex] at (V1) {}; \node[vertex] at (V2) {};
\draw[muF] (-0.4,-1.75) rectangle (1.0,1.75);
\node[lab,text=blue!60!black,anchor=south] at (0.3,1.8) {$\muF$: what counts as\\``inside the proton''};
\draw[muR] (2.1,-0.55) rectangle (5.1,0.55);
\node[lab,text=orange!80!black,anchor=north] at (3.6,-0.6) {$\muR$: at which scale\\$\alphas$ is evaluated, twice};
\node[lab,text=black!65,text width=6.9cm,anchor=north west] at (-0.4,-2.35) {Radiation harder than $\muF$ belongs to $\hat\sigma$, softer radiation to the PDFs; the split is a convention and the full answer does not depend on it, but a leading-order answer does. \PYTHIA sets both to $\muR^2=\muF^2=\pthat^2+\mb^2$ (scale codes 1--3).};
% ---------------- (b) the ladder of scales ---------------------------------------------------------
\begin{scope}[xshift=9.6cm]
\node[head,anchor=west] at (-0.3,3.0) {(b) the scales of our event, $\log_{10}(\mu/\mathrm{GeV})$};
% axis from log10 = -0.7 (0.2 GeV) to 4.4 (25 TeV): 1 decade = 2.6 cm; x = (log10 v + 0.7) * 2.6
\draw[ladder] (0,0) -- (13.4,0);
\foreach \l/\t in {0/{$1$},1/{$10$},2/{$10^2$},3/{$10^3$},4/{$10^4$}} {\draw[tick] ({(\l+0.7)*2.6},-0.1) -- ({(\l+0.7)*2.6},0.1); \node[lab,below] at ({(\l+0.7)*2.6},-0.1) {\t};}
\draw[thick,gray!70] (0.46,0) -- (0.46,0.6);  \node[lab,anchor=south west,text=gray!80!black] at (0.3,0.6) {$\Lambda_{\mathrm{QCD}}\approx0.3\GeV$: perturbation theory fails};
\draw[thick,gray!70] (3.59,0) -- (3.59,1.4);  \node[lab,above,text=gray!80!black] at (3.59,1.4) {$\mb=4.8\GeV$\\threshold, $\rho=4\mb^2/\shat$};
\draw[thick] (6.92,0) -- (6.92,2.3); \node[lab,above] at (6.92,2.3) {$m_Z=91\GeV$: $\alphas=0.1180$, the input};
\draw[thick,orange!80!black] (7.16,0) -- (7.16,-0.8); \node[lab,below,text=orange!80!black] at (7.16,-0.8) {$\muR=\muF=Q=113\GeV$\\$\alphas(Q^2)=0.1143$};
\draw[thick,orange!60!black,dashed] (6.38,0) -- (6.38,-1.9); \node[lab,anchor=north east,text=orange!60!black] at (6.45,-1.9) {$Q/2=57\GeV$: $\alphas=0.1272$\\$\alphas^2$ up by $24\%$};
\draw[thick,orange!60!black,dashed] (7.94,0) -- (7.94,-1.9); \node[lab,anchor=north west,text=orange!60!black] at (7.87,-1.9) {$2Q=227\GeV$: $\alphas=0.1042$\\$\alphas^2$ down by $17\%$};
\draw[thick,red!70!black] (8.57,0) -- (8.57,1.0); \node[lab,anchor=south west,text=red!70!black] at (8.45,1.0) {$\sqrt{\shat}=396\GeV$: $\alphas=0.0977$\\(what scale code 4 would use)};
\draw[thick,gray!70] (12.57,0) -- (12.57,0.4); \node[lab,anchor=south east,text=gray!80!black] at (12.7,0.4) {$\sqrt{s}=13.6\TeV$};
\node[lab,text=black!65,text width=12.6cm,anchor=north west] at (-0.3,-3.1) {The hard process lives one decade above $\mb$ and four above $\Lambda_{\mathrm{QCD}}$. Moving $\muR$ by a factor two moves $\alphas^2$, and with it the whole leading-order rate, by about $20\%$: that is the uncertainty of \cref{sec:ancestral_home:scales}, and no amount of Monte Carlo statistics reduces it.};
\end{scope}
\end{tikzpicture}
}
\caption[Where the two scales enter, and the scales of our event on a logarithmic ladder]{The scales of the hard process. (a) Where $\muF$ and $\muR$ enter the factorised picture: $\muF$ decides which radiation is counted as part of the proton and which as part of $\hat\sigma$, $\muR$ is the argument of the two powers of $\alphas$. (b) The scales of our event on a logarithmic ladder, with the coupling at each from section~[2] of the \code{-v} log and from the embedded table: the hard process lives one decade above $\mb$ and four above $\Lambda_{\mathrm{QCD}}$, and a factor of two in $\muR$ moves $\alphas^2$ by about $20\%$.}
\label{fig:ancestral_home:scales_sketch}
\end{figure}

Every leading-order calculation has to choose $\muR$ and $\muF$. \PYTHIA's
default for a $2\to2$ process with two equal-mass final particles is the
transverse mass,
\begin{equation}
  \muR^2 = \muF^2 = \pthat^2 + \mb^2 ,
  \label{eq:ancestral_home:scale_choice}
\end{equation}
which its manual calls scale codes~1--3 (they coincide for equal masses);
code~4 uses $\shat$ and code~5 a fixed value. The multipliers
\code{SigmaProcess:renormMultFac} and \code{factorMultFac} multiply the
\emph{squared} scale, a source of factor-of-two confusion when one tries to
``double the scale''.

\paragraph{The running coupling.} What the choice of $\muR$ changes is the
value of $\alphas$, through the renormalisation-group equation
$\mathrm{d}\alphas/\mathrm{d}\ln\mu^2=-\beta_0\alphas^2/(4\pi)-\beta_1\alphas^3/(4\pi)^2-\dots$,
whose two-loop solution is the function \PYTHIA evaluates:
\begin{equation}
  \alphas(\mu^2)=\frac{4\pi}{\beta_0L}\left[1-\frac{\beta_1}{\beta_0^2}\frac{\ln L}{L}\right],
  \qquad L=\ln\frac{\mu^2}{\Lambda^2},\quad
  \beta_0=11-\tfrac23n_f,\quad \beta_1=102-\tfrac{38}{3}n_f ,
  \label{eq:ancestral_home:running}
\end{equation}
with $n_f=5$ active flavours above the $b$ threshold, so $\beta_0=23/3$ and
$\beta_1=116/3$. The one free parameter $\Lambda$ is fixed by the input
$\alphas(m_Z^2)=0.118$; for \cref{eq:ancestral_home:running} with five
flavours that gives $\Lambda_5=0.226\GeV$, and the formula then reproduces
the embedded table to five digits at the scale of our event,
$\alphas(12\,842\GeV^2)=0.11425$, and at a quarter of it, $0.12715$. Above
the top threshold \PYTHIA switches to $n_f=6$ and the simple formula starts
to drift, by $0.4\%$ at $4Q^2$ and $1\%$ at $\shat$; the table carries the
exact values. The numbers that matter for the next section are the two
ratios the log prints: $\alphas^2(Q^2/4)/\alphas^2(Q^2)=1.239$ and
$\alphas^2(4Q^2)/\alphas^2(Q^2)=0.831$. A factor of two in the scale is
$24\%$ one way and $17\%$ the other in the leading-order rate, before the
PDFs have moved at all.

The coupling in the hard process is a separate object from the coupling in
the parton shower and from the one stored with the PDF set. The CP5
tune~\cite{CMS:2019csb} sets \code{SigmaProcess:alphaSvalue = 0.118} with
second-order running (\code{alphaSorder = 2}) for the matrix element, and the
same for the initial- and final-state showers and for multi-parton
interactions. Using an NNLO PDF set does not make the matrix element NNLO,
and in general using the PDF set's own $\alphas$ would not reproduce
\PYTHIA. At our event the two happen to coincide,
$\alphas(12\,842\GeV^2)=0.1143$ either way
(\cref{tab:ancestral_home:pdf_values}), because both are second-order
runnings from the same $\alphas(m_Z)$ and the scale is close to $m_Z$;
at the $30\GeV$ draft's $Q^2=1178\GeV^2$ they differed in the third digit.


\begin{figure}[htbp]\centering
\includegraphics[page=6,width=\textwidth]{ME/Standalone/theory_qqbar_bbbar.pdf}
\caption[The running coupling of the embedded table and the scales the integral visits]{The coupling and the scale, from \file{plot_theory_ME.py}. (a) \PYTHIA's second-order $\alphas(Q^2)$ as embedded in the standalone script, with the values at $m_Z^2$, at our event's $Q^2=12\,842\GeV^2$ and at a quarter and four times that scale. (b) The distribution of the cross section in the scale $Q=\sqrt{\pthat^2+\mb^2}$ that \cref{eq:ancestral_home:scale_choice} assigns to each point: half of the rate is evaluated below $Q=122\GeV$, within a quarter of the cut, so the single number $\alphas(113\GeV)$ of our event is representative of the whole sample.}
\label{fig:ancestral_home:scales_plot}
\end{figure}

\section{From formula to number: the cross section step by step}
\label{sec:ancestral_home:from_formula_to_number}

We now have everything needed to turn \cref{eq:ancestral_home:factorisation}
into $345\pb$, and I want you to see the whole machine, not just the output.
The steps below are, one for one, the steps of the standalone script you
ran in \cref{sec:ancestral_home:standalone_code}; the function names in
brackets point to it.

\begin{figure}[htbp]\centering
\resizebox{\textwidth}{!}{% ---------------------------------------------------------------------------
% The seven steps of the standalone integrator as a flow chart, with the
% numbers of our event's point in the unit cube (section [2] of the -v log of
% standalone_qqbar_bbbar_xsec.py) on the left and the pass of 2^16 points
% (section [3]) on the right. Function names in brackets.
% ---------------------------------------------------------------------------
\begin{tikzpicture}[x=1cm,y=1cm,
  every node/.style={font=\footnotesize},
  box/.style={draw,rounded corners=3pt,thick,fill=gray!8,align=left,inner sep=5pt,minimum height=11mm,text width=5.4cm},
  stepbox/.style={box,fill=orange!12},
  numbox/.style={draw,rounded corners=3pt,thin,fill=blue!6,align=left,inner sep=4pt,text width=8.2cm,font=\scriptsize},
  stepno/.style={draw,circle,inner sep=1pt,font=\scriptsize\bfseries,fill=orange!40,minimum size=5mm},
  flow/.style={-{Stealth[length=2.5mm,width=2mm]},thick},
  fn/.style={font=\scriptsize\ttfamily,text=black!65},
  head/.style={font=\small\bfseries}]
\node[head,anchor=west] at (-0.6,1.5) {the recipe, one Sobol point at a time};
\node[head,anchor=west] at (6.5,1.5) {the same steps for our event's point $(u_1,u_2,u_3)=(0.1617,0.1331,0.9505)$};
\node[stepbox] (S1) at (2.4,0) {\textbf{1\; choose the variables}\\ $(u_1,u_2,u_3)\mapsto(\ln\tau,\,Y,\,|\cos\theta|)$, each map linear on its interval};
\node[stepbox,below=8mm of S1] (S2) {\textbf{2\; compute the kinematics}\\ $\shat=\tau s$, $x_{1,2}$, $p^*$, $\pthat$, $Q^2=\pthat^2+\mb^2$};
\node[stepbox,below=8mm of S2] (S3) {\textbf{3\; evaluate the ingredients}\\ $\alphas(Q^2)$, ten $xf$ values, $\mathrm{d}\hat\sigma/\mathrm{d}\cos\theta$};
\node[stepbox,below=8mm of S3] (S4) {\textbf{4\; form the weight}\\ Jacobian $\times$ matrix element $\times$ PDF products of both orientations};
\node[stepbox,below=8mm of S4] (S5) {\textbf{5\; average over the cube}\\ $2^{16}$ scrambled Sobol points; $\sigma=\langle w\rangle$};
\node[stepbox,below=8mm of S5] (S6) {\textbf{6\; estimate the error}\\ $R$ scrambles, spread of the $R$ averages};
\node[stepbox,below=8mm of S6] (S7) {\textbf{7\; convert the units}\\ $\mathrm{GeV}^{-2}\to$ pb};
\foreach \a/\b in {S1/S2,S2/S3,S3/S4,S4/S5,S5/S6,S6/S7} {\draw[flow] (\a) -- (\b);}
\foreach \n/\s in {1/S1,2/S2,3/S3,4/S4,5/S5,6/S6,7/S7} {\node[stepno] at ($(\s.west)+(-0.35,0)$) {\n};}
\node[fn,anchor=south east] at (S1.north east) {born\_weights (steps 1--4)};
\node[fn,anchor=south east] at (S5.north east) {integrate (steps 5--6)};
% the numbers
\node[numbox,anchor=west] (N1) at (6.5,0) {$\ln\tau=\ln\tau_{\min}+u_1\ln(1/\tau_{\min})=-8.437+0.1617\times8.437=-7.072$\\
  $Y=(2u_2-1)\,Y_{\max}=(-0.734)\times3.536=-2.595$\\
  $|\cos\theta|=u_3\,c_{\max}=0.9505\times0.8631=0.820$};
\node[numbox,anchor=west] (N2) at (6.5,0 |- S2) {$\tau=8.483\times10^{-4}$, $\shat=156\,895\GeV^2$, $\sqrt{\shat}=396.1\GeV$\\
  $x_1=\sqrt\tau e^{Y}=2.174\times10^{-3}$, $x_2=\sqrt\tau e^{-Y}=0.3901$\\
  $p^*=198.0\GeV$, $\pthat=p^*\sin\theta=113.2\GeV$, $Q^2=12\,842\GeV^2$};
\node[numbox,anchor=west] (N3) at (6.5,0 |- S3) {$\alphas(Q^2)=0.11425$ \quad[TableAlphaS]\\
  $x_1f_{\bar u}=1.1455$, $x_2f_u=0.2107$, $x_1f_u=1.2566$, $x_2f_{\bar u}=0.0012$, \dots\quad[TablePDF.xf]\\
  $\mathrm{d}\hat\sigma/\mathrm{d}\cos\theta=18.92\pb$ \quad[dsigma\_dcostheta]};
\node[numbox,anchor=west] (N4) at (6.5,0 |- S4) {Jacobian $8.437\times7.072\times1.726=103.0$\\
  $w_u=103.0\times18.92\times(1.2566\times0.0012+1.1455\times0.2107)=473.2\pb$\\
  $w_d=155.3$, $w_s=12.3$, $w_c=9.1$, $w_b=0.9$;\quad $\textstyle\sum_q w_q=650.7\pb$};
\node[numbox,anchor=west] (N5) at (6.5,0 |- S5) {running average: $394.5$ ($N=16$), $338.7$ ($2^8$), $345.22$ ($2^{10}$), $345.28$ ($2^{14}$), $345.320\pb$ ($2^{16}$)\\
  per flavour: $u\bar u$ $178.3$, $d\bar d$ $118.6$, $s\bar s$ $30.3$, $c\bar c$ $12.7$, $b\bar b$ $5.5\pb$};
\node[numbox,anchor=west] (N6) at (6.5,0 |- S6) {eight scrambles: $345.320$, $.318$, $.324$, $.321$, $.323$, $.321$, $.324$, $.320\pb$\\
  $\bar\sigma=345.322\pb$, $\Delta\sigma=0.001\pb$ ($3\times10^{-6}$)};
\node[numbox,anchor=west] (N7) at (6.5,0 |- S7) {$3.893\,793\,72\times10^{8}\pb\,\mathrm{GeV}^2$, applied inside \texttt{dsigma\_dcostheta}, so every weight is already in pb};
\foreach \s/\n in {S1/N1,S2/N2,S3/N3,S4/N4,S5/N5,S6/N6,S7/N7} {\draw[flow,gray!60] (\s.east) -- (\n.west);}
\end{tikzpicture}
}
\caption[The seven steps of the integrator, with the numbers of our event's point]{The seven steps of the integrator as a flow chart. Left: the recipe, one Sobol point at a time; steps~1 to~4 are the function \code{born\_weights}, steps~5 and~6 the function \code{integrate}. Right: the same steps for the point of the unit cube that corresponds to our event, with every intermediate number from sections~[2] and~[3] of the \code{-v} log, ending in the weight $650.7\pb$ of that one point and the average $345.320\pb$ over the $65\,536$ points of the first pass.}
\label{fig:ancestral_home:recipe_flow}
\end{figure}

\paragraph{Step 1: choose the integration variables.} Every Monte Carlo
integrator wants to sample a unit cube, so we map
$(u_1,u_2,u_3)\in[0,1]^3$ onto $\ln\tau$, $Y$ and $|\cos\theta|$, each map
being linear on its interval:
\begin{align}
  \ln\tau &= \ln\tau_{\min} + u_1\,\ln\frac{1}{\tau_{\min}},
  & \mathrm{d}\ln\tau &= \ln\frac{1}{\tau_{\min}}\;\mathrm{d}u_1, \notag\\
  Y &= (2u_2-1)\,Y_{\max}(\tau),
  & \mathrm{d}Y &= 2Y_{\max}\;\mathrm{d}u_2, \label{eq:ancestral_home:maps}\\
  |\cos\theta| &= u_3\,c_{\max}(\tau),
  & \mathrm{d}|\cos\theta| &= c_{\max}\;\mathrm{d}u_3 . \notag
\end{align}
Three remarks, one per line. Sampling $\ln\tau$ rather than $\tau$ is not
cosmetic: partonic energies at the LHC span orders of magnitude, and the
PDFs favour small $x$ so strongly that a uniform sample in $\tau$ would waste
almost all its points near $\tau=1$ where nothing happens. The rapidity
interval depends on $\tau$ through \cref{eq:ancestral_home:ymax}, so the
cube maps onto a triangle, not a rectangle. And only positive $\cos\theta$ is
sampled, with a factor~2 to follow, because the matrix element and all our
cuts are even in $\cos\theta$ [\code{born\_weights}].

\paragraph{Step 2: compute the kinematics.} $\shat=\tau s$, $x_{1,2}$ from
\cref{eq:ancestral_home:tauY}, $p^{*}$ from \cref{eq:ancestral_home:pstar},
$\pthat^2=p^{*2}(1-\cos^2\theta)$, and the scale $Q^2=\pthat^2+\mb^2$. A
point with $p^{*}<\pthat^{\min}$ gets weight zero.

\paragraph{Step 3: evaluate the ingredients at that point.} $\alphas(Q^2)$
[\code{TableAlphaS}], the ten PDF values $x_1f_q$, $x_1f_{\bar q}$,
$x_2f_q$, $x_2f_{\bar q}$ for $q=d,u,s,c,b$ [\code{TablePDF.xf}], and the
matrix element $\mathrm{d}\hat\sigma/\mathrm{d}\cos\theta$
[\code{dsigma\_dcostheta}].

\paragraph{Step 4: form the weight.} Put \cref{eq:ancestral_home:xf} and
\cref{eq:ancestral_home:maps} into \cref{eq:ancestral_home:factorisation}.
The integrand times the volume of the mapped region is, for each incoming
flavour $q$,
\begin{equation}
  w_q(u_1,u_2,u_3)=
  \underbrace{\ln\frac{1}{\tau_{\min}}}_{u_1\to\ln\tau}\;
  \underbrace{2Y_{\max}}_{u_2\to Y}\;
  \underbrace{2c_{\max}}_{u_3\to\pm\cos\theta}\;
  \frac{\mathrm{d}\hat\sigma}{\mathrm{d}\cos\theta}\;
  \bigl[\,x_1f_q(x_1)\,x_2f_{\bar q}(x_2)+x_1f_{\bar q}(x_1)\,x_2f_q(x_2)\,\bigr],
  \label{eq:ancestral_home:weight}
\end{equation}
where the bracket holds both beam orientations and the second factor of~2 in
front of $c_{\max}$ restores the negative hemisphere. Because the cube has
unit volume, the cross section is simply the \emph{average} of the weight
over the cube,
\begin{equation}
  \sigma_q=\int_0^1\!\!\int_0^1\!\!\int_0^1 w_q\,\mathrm{d}u_1\mathrm{d}u_2\mathrm{d}u_3
  \approx\frac1N\sum_{i=1}^{N}w_q(\mathbf u_i),\qquad
  \sigma=\sum_{q=d,u,s,c,b}\sigma_q .
  \label{eq:ancestral_home:average}
\end{equation}
There is nothing else. Every Monte Carlo cross section you will ever read
out of a generator log is an average of weights of this kind.

\paragraph{Step 5: average over many points.} Which points? Not
pseudo-random ones. Independent random points cluster and leave gaps, and
their error falls only as $1/\sqrt N$. A Sobol sequence is a
\emph{quasi-random}, low-discrepancy set that fills the cube deliberately
evenly; for a smooth three-dimensional integrand like ours it converges much
faster. The scripts use $2^{16}=65\,536$ points per pass (powers of two
preserve the sequence's balance), and a single pass already gives four
significant digits [\code{integrate}].

\paragraph{Step 6: estimate the numerical error.} A plain Sobol sequence is
deterministic, so running it twice teaches nothing. \emph{Scrambling}
randomises the sequence while keeping its even coverage, and $R$ passes with
different scrambles give $R$ nearly independent estimates
$\sigma^{(r)}$. The result and its error are
\begin{equation}
  \bar\sigma=\frac1R\sum_{r=1}^{R}\sigma^{(r)},\qquad
  \Delta\sigma=\frac{1}{\sqrt R}\sqrt{\frac{1}{R-1}\sum_{r=1}^{R}\bigl(\sigma^{(r)}-\bar\sigma\bigr)^2}.
  \label{eq:ancestral_home:error}
\end{equation}
The flavour terms are correlated (they share the same points), so the total
is summed within each pass before the spread is taken.

\paragraph{Step 7: convert units.} Multiply by $3.893\,793\,72\times10^{8}$
to go from $\GeV^{-2}$ to pb.

\paragraph{The numbers.} Section~[3] of the log is the first Sobol pass
told in full. Of its $65\,536$ points all but $31$ have a positive weight;
the weights range from zero to $4576\pb$, their median is $1.0\pb$ and their
mean is $345.32\pb$, which is the cross section. The mean is three hundred
times the median because the integrand is steep: the largest tenth of the
weights carries $75\%$ of the integral and the largest hundredth $11\%$, and
in terms of $u_1$ the first tenth of the interval, $\sqrt{\shat}$ between
$200$ and $305\GeV$, holds $63\%$ of it, the top seventy per cent of the
interval, above $710\GeV$, under two. Our event's own point weighs
$650.7\pb$, about twice the mean. The running average swings as predicted:
$26\pb$ after four points, $401$ after eight, $394$ after sixteen, $339$
after $256$, $345.22$ after $1024$, and from $2^{14}$ points on it does not
move beyond the second decimal, $345.28\to345.319\to345.320\pb$. The
relative spread of the weights is $2.37$, so pseudo-random points would give
an error of $2.37\times345/\sqrt{65\,536}=3.19\pb$ for the same effort;
eight scrambled Sobol passes give $0.001\pb$
(\cref{tab:ancestral_home:three_numbers}). \Cref{fig:ancestral_home:convergence_plot}
shows the weight distribution, the running average and the error against
$N$ for both kinds of points.

\begin{figure}[htbp]\centering
\includegraphics[page=7,width=\textwidth]{ME/Standalone/theory_qqbar_bbbar.pdf}
\caption[The weights, the running average and the error of the Sobol integration]{From formula to number, from \file{plot_theory_ME.py}. (a) The distribution of the $65\,536$ weights of the first Sobol pass on a logarithmic scale; the mean, which is the cross section, and the weight of our event's point are marked. (b) The running average against the number of points for the scrambled Sobol sequence of the script and for pseudo-random points with the same seed; the band is the \PYTHIA result of \cref{sec:ancestral_home:cmssw} with its statistical error. (c) The spread of eight independent estimates against $N$: pseudo-random points follow $N^{-1/2}$, the Sobol sequence is close to $N^{-1}$ and is a thousand times more precise at $2^{16}$ points.}
\label{fig:ancestral_home:convergence_plot}
\end{figure}

The results are in \cref{tab:ancestral_home:three_numbers} and, per flavour,
in \cref{tab:ancestral_home:flavours}. Up quarks dominate, more so than at
$30\GeV$, because the higher cut pushes the valence fraction $x_2$ up to
where $u$ is twice as abundant as $d$. The dependence on the \pt cut
(\cref{tab:ancestral_home:ptscan}) is steep: raising $\pthat^{\min}$ from
$30$ to $100\GeV$ costs a factor~54, and from $100$ to $200\GeV$ a further
factor~13.

\begin{table}[htbp]\centering
\caption[Contribution of each incoming flavour]{Contribution of each
incoming flavour to the default cross section, from the full script
(eight repeats of $2^{16}$ Sobol points, \file{result.json} in
\file{example/ME/Standalone/}).}
\label{tab:ancestral_home:flavours}
\begin{tabular}{lrr}\toprule
Incoming pair & $\sigma$ [pb] & fraction \\ \midrule
$u\bar u$ & $178.316$ & $51.6\%$ \\
$d\bar d$ & $118.568$ & $34.3\%$ \\
$s\bar s$ & $30.269$ & $8.8\%$ \\
$c\bar c$ & $12.697$ & $3.7\%$ \\
$b\bar b$ & $5.474$ & $1.6\%$ \\ \midrule
total & $345.324 \pm 0.001$ & \\
\bottomrule
\end{tabular}
\end{table}

\begin{table}[htbp]\centering
\caption[Dependence on the generator-level cut $\pthat^{\min}$]{Dependence
on the generator-level cut $\pthat^{\min}$, all other settings at their
defaults, from single passes of the standalone script with
\code{--ptmin}. The \CMSSW and \MG references of
\cref{tab:ancestral_home:three_numbers} apply to the $100\GeV$ row only.}
\label{tab:ancestral_home:ptscan}
\begin{tabular}{rr}\toprule
$\pthat^{\min}$ [GeV] & $\sigma$ [pb] \\ \midrule
30 & $18\,506$ \\
50 & $3\,650.5$ \\
100 & $345.32$ \\
150 & $79.74$ \\
200 & $26.96$ \\
300 & $5.350$ \\
\bottomrule
\end{tabular}
\end{table}

\subsection{What this number is, and what it is not}
\label{sec:ancestral_home:caveats}

The agreement in \cref{tab:ancestral_home:three_numbers} was obtained without
fitting or adjusting anything to the \PYTHIA result. It holds because the
integrator reproduces \PYTHIA's conventions, and each convention is a caveat
you should carry with you to the later stops:
\begin{itemize}
  \item It is the \emph{Born} hard-process rate, unweighted and unsuppressed.
    Nothing has showered, hadronised or been clustered. It is not a $b$-jet
    cross section, and a cut on $\pthat$ is not a cut on jet \pt.
  \item The $|y|$ or $|\eta|$ cuts the full integrator offers act on the two
    Born quarks, not on reconstructed jets.
  \item The incoming-$b$ term uses the $s$-channel formula only
    (\cref{sec:ancestral_home:gg_and_flavours}); the $gg$ channel is absent by
    choice.
  \item PDFs are frozen at grid edges, clipped at zero, and the heavy sea is
    symmetrised (\cref{sec:ancestral_home:pdfs}).
  \item The coupling is \PYTHIA's second-order $\alphas$ with
    $\alphas(m_Z)=0.118$, not the PDF set's.
  \item The quoted errors are integration or sampling errors. They say nothing
    about the accuracy of leading-order QCD, which is the subject of
    \cref{sec:ancestral_home:scales}.
  \item Points in a \pt scan share the same Sobol seeds and are therefore not
    statistically independent bins.
  \item User hooks, process weights, phase-space biasing, nuclear PDFs and
    any post-shower selection that a production configuration might add are
    not reproduced.
\end{itemize}

\section{Scale choices and their uncertainties}
\label{sec:ancestral_home:scales}

I have been quoting cross sections to six digits, and it is time to confess
how many of them mean anything. A leading-order cross section depends on
$\muR$ through $\alphas^2$ and on $\muF$ through the PDFs, and neither
dependence is compensated by anything at this order. The conventional
estimate of the missing higher orders is to vary both scales by factors of
two around the central choice \cref{eq:ancestral_home:scale_choice}.
\Cref{tab:ancestral_home:variations} shows the result for our process,
computed with the full integrator, and \cref{fig:ancestral_home:seven_point_sketch}
arranges the seven values that enter the standard prescription on the
$(\muR,\muF)$ grid. Guess first whether a factor of two in $\muR$ or in
$\muF$ matters more, and in which direction each moves the rate.

\begin{figure}[htbp]\centering
\resizebox{\textwidth}{!}{% ---------------------------------------------------------------------------
% The seven-point scale variation as a grid in (mu_R, mu_F), with the cross
% sections of table tab:ancestral_home:variations (qqbar_bbbar_xsec.py through
% scan_scales.py). The two opposite-extreme corners are not part of the
% prescription.
% ---------------------------------------------------------------------------
\begin{tikzpicture}[x=1cm,y=1cm,
  every node/.style={font=\footnotesize},
  cell/.style={draw,thick,minimum size=24mm,align=center,font=\footnotesize},
  up/.style={cell,fill=red!18},
  down/.style={cell,fill=blue!18},
  mid/.style={cell,fill=gray!10},
  skip/.style={cell,fill=gray!4,pattern=north east lines,pattern color=gray!40,text=black!60,font=\scriptsize},
  lab/.style={font=\scriptsize,align=center},
  head/.style={font=\small\bfseries},
  axis/.style={-{Stealth},thick}]
\node[head,anchor=west] at (-1.9,4.6) {the seven-point variation: $\muR,\muF\in\{\tfrac12,1,2\}\times\mu_0$ without the two opposite extremes};
% grid: columns mu_R/2, mu_R, 2mu_R ; rows 2mu_F (top), mu_F, mu_F/2 (bottom)
\node[up]   (c00) at (0,2.5)  {$406.6\pb$\\$+17.7\%$};
\node[up]   (c10) at (2.5,2.5){$357.2\pb$\\$+3.4\%$};
\node[skip] (c20) at (5,2.5)  {not used\\($\muR\times2$, $\muF\times\tfrac12$\\opposite extremes)};
\node[up]   (c01) at (0,0)    {$425.3\pb$\\$+23.2\%$};
\node[mid]  (c11) at (2.5,0)  {\textbf{345.3 pb}\\central};
\node[down] (c21) at (5,0)    {$288.8\pb$\\$-16.4\%$};
\node[skip] (c02) at (0,-2.5) {not used\\($\muR\times\tfrac12$, $\muF\times2$\\opposite extremes)};
\node[down] (c12) at (2.5,-2.5){$330.2\pb$\\$-4.4\%$};
\node[down] (c22) at (5,-2.5) {$298.7\pb$\\$-13.5\%$};
\draw[axis] (-1.5,-3.9) -- (6.3,-3.9) node[below left,lab] {$\muR$};
\foreach \x/\t in {0/{$\muR/2$},2.5/{$\muR$},5/{$2\muR$}} {\node[lab,below] at (\x,-3.95) {\t};}
\draw[axis] (-1.5,-3.9) -- (-1.5,3.9) node[below left,lab] {$\muF$};
\foreach \y/\t in {2.5/{$2\muF$},0/{$\muF$},-2.5/{$\muF/2$}} {\node[lab,left] at (-1.55,\y) {\t};}
% the envelope
\node[lab,text=black!65,text width=7.2cm,anchor=north west] at (7.0,3.6) {\textbf{Reading the grid.} Red cells raise the cross section, blue cells lower it. Moving along a row ($\muR$) changes $\alphas^2$ and is the large effect; moving along a column ($\muF$) changes the PDFs and is the small one, because the rising sea at $x_1$ and the falling valence at $x_2$ nearly cancel at our event.\\[4pt]
  \textbf{The envelope} is the smallest and largest of the seven values: $288.8$ to $425.3\pb$, that is\\[2pt]
  \mbox{}\hfill $\sigma_{\mathrm{LO}}=345.3\,^{+80.0}_{-56.5}\pb=345.3\,^{+23\%}_{-16\%}\pb$.\hfill\mbox{}\\[4pt]
  The \CMSSW statistical error ($\pm6.1\pb$) and the integration error ($\pm0.001\pb$) would be invisible on this scale.};
\end{tikzpicture}
}
\caption[The seven-point scale variation on the $(\muR,\muF)$ grid]{The seven-point scale variation. The nine combinations of $\muR,\muF\in\{\tfrac12,1,2\}\times\mu_0$ with $\mu_0^2=\pthat^2+\mb^2$, without the two opposite extremes, with the cross sections of \cref{tab:ancestral_home:variations}. Red cells raise the rate, blue cells lower it; the envelope of the seven values is the quoted theoretical uncertainty.}
\label{fig:ancestral_home:seven_point_sketch}
\end{figure}

\begin{table}[htbp]\centering
\caption[Scale, flavour, mass and PDF-convention variations]{Variations of
the default calculation, computed with \file{qqbar_bbbar_xsec.py} ($2^{15}$
Sobol points, two repeats; integration errors below $0.002\%$). The central
value is $345.32\pb$.}
\label{tab:ancestral_home:variations}
\begin{tabular}{llrr}\toprule
Variation & Option & $\sigma$ [pb] & change \\ \midrule
$\muR\times2$              & \code{--mur-factor 2}                   & $288.83$ & $-16.4\%$ \\
$\muR\times\tfrac12$       & \code{--mur-factor 0.5}                 & $425.34$ & $+23.2\%$ \\
$\muF\times2$              & \code{--muf-factor 2}                   & $357.19$ & $+3.4\%$ \\
$\muF\times\tfrac12$       & \code{--muf-factor 0.5}                 & $330.19$ & $-4.4\%$ \\
$\muR,\muF\times2$         & \code{--mur-factor 2 --muf-factor 2}     & $298.70$ & $-13.5\%$ \\
$\muR,\muF\times\tfrac12$  & \code{--mur-factor 0.5 --muf-factor 0.5} & $406.61$ & $+17.7\%$ \\ \midrule
no incoming $b$            & \code{--nquark 4}                       & $339.85$ & $-1.6\%$ \\
$\mb=4.7\GeV$              & \code{--mb 4.7}                         & $345.37$ & $+0.01\%$ \\
negative PDFs kept         & \code{--signed-pdfs}                    & $345.04$ & $-0.08\%$ \\
\bottomrule
\end{tabular}
\end{table}

\paragraph{The estimate you can do by hand.} The size of the $\muR$ rows
follows from \cref{eq:ancestral_home:running} alone. Since $\sigma\propto\alphas^2(\muR^2)$
at this order, and since to one loop
$\mathrm{d}\alphas/\mathrm{d}\ln\mu^2=-\beta_0\alphas^2/(4\pi)$,
\begin{equation}
  \frac{\Delta\sigma}{\sigma}\simeq2\,\frac{\Delta\alphas}{\alphas}
  \simeq-\frac{\beta_0\,\alphas(\mu_0^2)}{2\pi}\,\ln\frac{\muR^2}{\mu_0^2}
  =\mp\frac{23}{3}\cdot\frac{0.1143}{2\pi}\cdot\ln4=\mp19\%
  \quad\text{for }\muR=2\mu_0,\ \tfrac12\mu_0 .
  \label{eq:ancestral_home:scale_estimate}
\end{equation}
The full integrator gives $-16.4\%$ and $+23.2\%$: the two-loop running is
slower than the one-loop estimate above $\mu_0$ and faster below it, which
is what the exact ratios $0.831$ and $1.239$ of the log already said, and the
rest is the spread of $Q$ over the sample (\cref{fig:ancestral_home:scales_plot}).
The $\muF$ rows cannot be estimated this simply because they are the
difference of two DGLAP flows; at our event the log gives $+0.6\%$ and
$-1.1\%$ for the product of the two densities, and the integrated result,
$+3.4\%$ and $-4.4\%$, is of the same small size and the same sign. The
envelope of the seven points is then
\begin{equation}
  \sigma_{\mathrm{LO}}=345.3\,^{+80.0}_{-56.5}\pb
  =345.3\,^{+23\%}_{-16\%}\pb ,
  \label{eq:ancestral_home:envelope}
\end{equation}
and \cref{fig:ancestral_home:scale_variation_plot} shows how the three
curves behind it look between the extremes.

\begin{figure}[htbp]\centering
\includegraphics[page=8,width=\textwidth]{ME/Standalone/theory_qqbar_bbbar.pdf}
\caption[Scale variations of the leading-order cross section]{Scale variations of the leading-order cross section, computed with the full integrator through \file{scan_scales.py} and drawn by \file{plot_theory_ME.py}. (a) The cross section against the scale factor for $\muR$ alone, $\muF$ alone and both together, on a logarithmic axis in the factor; the band is the \PYTHIA result with its $1.8\%$ statistical error. (b) The same seven values on the $(\muR,\muF)$ grid, coloured by the change.}
\label{fig:ancestral_home:scale_variation_plot}
\end{figure}

Three lessons are in the table. First, the scale uncertainty of this
leading-order number is of order $20\%$, three orders of magnitude larger
than the integration errors and ten times the statistical error of the
\CMSSW run; agreement between two leading-order codes at $0.3\%$ says
nothing about agreement with nature. Second, $\muR$ and $\muF$ pull in
opposite directions: a larger $\muR$ lowers $\alphas$, while a larger $\muF$
raises the small-$x$ sea that our $\bar u$ came from, because DGLAP evolution
feeds the sea from the gluon as $Q^2$ grows, and at the same time lowers the
valence $u$ at $x=0.39$. At $\pthat\ge30\GeV$ the sea won clearly and the
$\muF$ variation was $\pm15\%$; at $100\GeV$ the two nearly cancel and it
is $\pm4\%$. Exercise~6 lets you watch this happen. Varying both scales
together therefore under-estimates the uncertainty, which is why the standard
prescription is the envelope of a seven-point variation ($\muR,\muF$ each
$\times\tfrac12,1,2$ without the two extreme opposite combinations). Third,
the ``convention'' rows are small: whether negative PDF values are clipped is
irrelevant, the $b$ mass is invisible at this \pt, and the incoming-$b$
approximation of \cref{sec:ancestral_home:gg_and_flavours} is a $1.6\%$
effect.

Beyond scales, a complete uncertainty would add the PDF uncertainty from the
other members of the set (100 replicas for NNPDF3.1, run with
\code{--member}), the variation of $\alphas(m_Z)$ by $\pm0.001$
(\code{--alphas-mz}), and the choice of $\mb$, which enters through $\rho$
and through the scale. The real cure is the next order: at NLO the $\muR$
and $\muF$ dependence of the Born term is cancelled by the explicit
logarithms of the one-loop and real-emission corrections, and the residual
uncertainty shrinks to about $10\%$. That is what \MG at
NLO~\cite{Alwall:2014hca} provides, and why CMS uses NLO samples where the
cross section itself matters. For following one jet through the detector,
leading order is enough; for a measurement it is not.

\section{Validation from official packages}
\label{sec:ancestral_home:validation}

A script you wrote yourself and a derivation you did yourself can be wrong
in the same way. The last word at every stop therefore goes to the official
packages: the code CMS actually runs, and an independent generator that
shares none of its source. Neither is a black box any more; after the
previous sections you know what every setting in them means, and the point
of this section is to read their configurations and logs and find our
$345\pb$ in them.

\subsection{How CMS does it in practice: the GEN step in \CMSSW}
\label{sec:ancestral_home:cmssw}

Now let us see how the professionals do it. In CMS the hard process is not
run by hand; it is one step of a production chain driven by
\code{cmsDriver.py}. The configuration in
\file{example/ME/CMSSW/step1_GEN_cfg.py} was generated with the command in
its header, here broken into lines:
\begin{minted}[linenos=false,frame=single,framesep=1.5mm,fontsize=\small,breaklines]{bash}
cmsDriver.py Configuration/GenProduction/python/BBbar_14TeV_TuneCP5_cfi.py \
  --conditions auto:phase1_2024_realistic --era Run3_2024 --geometry DB:Extended \
  -n 1 -s GEN --datatier GEN --eventcontent FEVTDEBUG \
  --fileout file:output_step1_GEN.root --python_filename step1_GEN_cfg.py \
  --no_exec --beamspot DBrealistic \
  --customise_commands "process.generator.comEnergy = cms.double(13600.0)"
\end{minted}
The pieces are: a \emph{fragment} (the \code{\_cfi.py} file, which carries the
\PYTHIA settings), the detector \emph{conditions} and \emph{era}, the
\emph{step} (\code{GEN} only), the output data tier and event content, and a
customisation that overrides the fragment's $14\TeV$ with the Run~3 energy
of $13.6\TeV$. To run it:
\begin{minted}[linenos=false,frame=single,framesep=1.5mm,fontsize=\small]{bash}
cmsrel CMSSW_15_0_5
cd CMSSW_15_0_5/src
cmsenv
cmsRun step1_GEN_cfg.py 2>&1 | tee step1_GEN_cfg.log
\end{minted}

\subsubsection{Reading the configuration}

Open the file on GitHub. Most of it is boilerplate that every \CMSSW job
carries. The physics is in the module \code{process.generator}, an
\code{EDFilter} of type \code{Pythia8ConcurrentGeneratorFilter}, and in
particular in three blocks of its \code{PythiaParameters}:
\begin{itemize}[nosep]
  \item \code{processParameters}: the three lines that define our process,
    \code{HardQCD:gg2bbbar = off}, \code{HardQCD:qqbar2bbbar = on} and
    \code{PhaseSpace:pTHatMin = 100.}. Everything in
    \cref{sec:ancestral_home:matrix_element} follows from these.
  \item \code{pythia8CP5Settings}: the CMS tune~\cite{CMS:2019csb}. It fixes
    the PDF through \code{PDF:pSet} to the NNPDF3.1 set of
    \cref{sec:ancestral_home:pdfs}, and sets the four couplings, for the matrix
    element, the two showers and the multi-parton interactions, to
    $\alphas(m_Z)=0.118$ at second order. It also carries the
    multi-parton-interaction and colour-reconnection parameters that later
    stops will meet.
  \item \code{pythia8CommonSettings}: production defaults such as the
    tolerance on energy--momentum conservation and the lifetime cut-off for
    particle decays.
\end{itemize}
Other lines worth noticing: \code{maxEvents.input = 1000} (the fragment
name says one event, the job was run with a thousand); the
\code{comEnergy} override at the end; the output module writing
\code{FEVTDEBUG} content to \file{output_step1_GEN.root}; and a commented-out
\code{customiseJitJet} line that hooks in the book's own analysis package,
used from the next chapter on.

\subsubsection{Reading the log}

\PYTHIA writes two things to the log that close the loop with
\cref{sec:ancestral_home:from_formula_to_number}. For the first event it
prints the ``Info Listing'': the two incoming partons with their $x$ and
PDF values, the partonic invariants $\shat$, $\hat t$, $\hat u$, $\pthat$
and $\hat\theta$, and the coupling at the scale $Q^2$. These are the
quantities of \cref{tab:ancestral_home:pdf_values} and
\cref{eq:ancestral_home:kinematics}; find them in your own log and check
them against the table. Immediately after it comes the hard-process listing
that we reproduce in \cref{sec:ancestral_home:our_jet}.

At the end of the job it prints the ``Event and Cross Section Statistics''
table: how many phase-space points were tried, how many selected and
accepted, and the cross section with its error, in mb. About three
phase-space points are tried for each unweighted event that is kept; the
cross section is the average weight of the tried points, not the acceptance
ratio, exactly as in \cref{eq:ancestral_home:average}. For our run the table
reads
\begin{equation*}
  \sigma(q\bar q\to\bbbar)=3.451\times10^{-7}\,\mathrm{mb}\pm6.078\times10^{-9}\,\mathrm{mb}
  =345.1\pm6.1\pb ,
\end{equation*}
$0.04$ standard deviations from our $345.32\pb$. The $1.8\%$ statistical
error of a thousand-event run is the largest number in
\cref{tab:ancestral_home:three_numbers}, and, as \cref{sec:ancestral_home:scales}
will show, still ten times smaller than the theoretical uncertainty. Nothing
here calls for a correction.

\subsection{The same calculation in \MG}
\label{sec:ancestral_home:madgraph}

A second generator is the best check of the first, because it shares none of
its code. \MG~\cite{Alwall:2014hca} builds the Feynman diagrams of a process
you type at a prompt, writes the matrix element as Fortran, integrates it
with its own adaptive sampler and writes the events in the Les Houches Event
(LHE) format. The recipe in \file{example/ME/MadGraph/qqbar_bbbar_xsec_MadGraph.md}
reproduces our number. Download and unpack the release, start the
interpreter, and type:
\begin{minted}[linenos=false,frame=single,framesep=1.5mm,fontsize=\small]{text}
MG5_aMC> define p = u c d s u~ c~ d~ s~
MG5_aMC> generate p p > b b~
MG5_aMC> output qq_to_bbar
MG5_aMC> launch
\end{minted}
The first line is the whole physics choice of this chapter in one command:
it redefines the proton as a bag of light quarks and antiquarks, without the
gluon, so that only $q\bar q\to\bbbar$ is generated. The second line asks
for the diagrams; with the default Standard Model, in which the $b$ is
massive and is not a constituent of the proton, there is exactly one, the
$s$-channel gluon of \cref{eq:ancestral_home:dsigma}. When \code{launch}
offers to edit the run card, set the beam energies, the PDF and the $b$
\pt cut:
\begin{minted}[linenos=false,frame=single,framesep=1.5mm,fontsize=\small]{text}
 6800.0   = ebeam1   ! beam 1 total energy in GeV
 6800.0   = ebeam2   ! beam 2 total energy in GeV
 nn23lo1   = pdlabel  ! PDF set
 303600   = lhaid    ! NNPDF31_nnlo_as_0118
 100.0    = ptb      ! minimum pT for the b quark
\end{minted}
and let the run proceed. The results summary of the run gives
\begin{equation*}
  \sigma(pp\to\bbbar)=344.3\pm1.0\pb\qquad(10\,000\text{ events}),
\end{equation*}
the last row of \cref{tab:ancestral_home:three_numbers}.

\paragraph{Why the agreement is good but not exact.} Do not read the
$0.3\%$ difference from \PYTHIA as a discrepancy, and do not read it as a
triumph either; several conventions differ, and they partly cancel:
\begin{itemize}[nosep]
  \item \textbf{Four incoming flavours, not five.} The \code{define p} line
    leaves the $b$ out of the proton, the 4FS road of
    \cref{sec:ancestral_home:gg_and_flavours}. The full script with
    \code{--nquark 4} gives $339.9\pb$, $1.6\%$ below the default.
  \item \textbf{The $b$ mass.} The default \code{sm} model uses
    $\mb=4.7\GeV$, \PYTHIA $4.8\GeV$. The script with \code{--mb 4.7} moves
    by $+0.01\%$: at $\pthat\ge100\GeV$ the mass is invisible.
  \item \textbf{The scale.} \MG's default dynamical scale for a $2\to2$
    process is built from the transverse masses of the final state after a
    $\kt$ clustering, not exactly \cref{eq:ancestral_home:scale_choice},
    and $\muR$ variations of a factor two move the result by $\pm20\%$
    (\cref{tab:ancestral_home:variations}).
  \item \textbf{The coupling.} With \code{pdlabel = nn23lo1}, \MG takes
    $\alphas$ from the PDF set; \PYTHIA runs its own. At this scale the two
    agree to four digits (\cref{tab:ancestral_home:pdf_values}), so this
    difference is negligible here but not in general.
\end{itemize}
The net effect of the first two is $-1.6\%$, and \MG comes out $1.3\%$
\emph{above} that, so its scale convention must be worth about a percent in
the other direction. Exercise~8 asks you to pin this down. The lesson is the
one of \cref{sec:ancestral_home:scales}: two leading-order codes agreeing to
a percent tells you that both were set up as intended, and nothing about
nature.

\paragraph{The LHE file.} \MG stores the result as weighted or unweighted
events in the LHE format, whose \code{<event>} blocks carry the same
information as \PYTHIA's hard-process listing of
\cref{sec:ancestral_home:our_jet}: one line
per particle with its PDG code, status ($-1$ incoming, $1$ outgoing), mother
indices, colour labels, four-momentum and mass, plus the event weight and the
values of $\alphas$ and the scale that were used. The plotting script of
\cref{sec:ancestral_home:our_jet} reads both formats. If you prefer to
follow an \MG event rather than a \PYTHIA one through the rest of the
journey, this is where you would branch off; the later stops do not care
who wrote the hard-process record.

\section{Our jet at this stage: the \texorpdfstring{$b$}{b}-quark four-vector in the hard-process record}
\label{sec:ancestral_home:our_jet}

Here is the mother at the moment of her birth. The file
\file{example/ME/event_ME.lhe} holds \PYTHIA's hard-process listing for the
event followed through the book, copied from the log of the \CMSSW job:

\filelisting[\widelisting]{text}{example/ME/event_ME.lhe}{the hard-process record of our event; momenta and masses in GeV}

Reading the columns: \code{status} $-12$ marks the two beam protons, $-21$ the
two incoming partons after they have been resolved out of the protons, and
$23$ the two outgoing partons of the hard process; negative status means the
particle has been ``used up'' and positive means it is still there for the
next stage. The colour columns show the colour flow: the $\bar u$ carries
anticolour~101, which is handed to the $\bar b$, and the $u$ carries
colour~102, handed to the $b$; the gluon in between never appears because it
is internal. The momentum sum at the bottom is that of the two incoming
partons, $p_z=14.785-2652.934=-2638.15\GeV$: the whole system recoils
backward, and its mass, $396.1\GeV$, is $\sqrt{\shat}$.

\begin{table}[htbp]\centering
\caption[Kinematics of the two quarks of our event]{Kinematics of the two
quarks in \file{event_ME.lhe}, computed by \file{example/ME/plot_journey_ME.py}.
The $b$ quark is the mother we follow.}
\label{tab:ancestral_home:kinematics}
\begin{tabular}{lrrrrrrrr}\toprule
 & $p_x$ & $p_y$ & $p_z$ & $E$ & $\pt$ & $\eta$ & $y$ & $\phi$ \\
 & [GeV] & [GeV] & [GeV] & [GeV] & [GeV] & & & [rad] \\ \midrule
$b$       & $-62.071$ & $-94.690$ & $-225.148$  & $252.059$  & $113.22$ & $-1.439$ & $-1.438$ & $-2.151$ \\
$\bar b$  & $62.071$  & $94.690$  & $-2413.001$ & $2415.661$ & $113.22$ & $-3.753$ & $-3.752$ & $0.991$ \\ \midrule
pair      & $0$ & $0$ & $-2638.149$ & $2667.719$ & $0$ & & $Y=-2.595$ & \\
\bottomrule
\end{tabular}
\end{table}

\Cref{tab:ancestral_home:kinematics} and \cref{fig:ancestral_home:journey}
show what the record means. The two quarks are exactly back to back in the
transverse plane ($\Delta\phi=\pi$) with $\pt=113.22\GeV$ each, just above
the $100\GeV$ generator cut, as most events in a steeply falling spectrum
are. The pair mass is $m_{\bbbar}=\sqrt{\shat}=396.1\GeV$, and
$\shat=x_1x_2s=2.174\times10^{-3}\times0.3901\times13600^2=156\,895\GeV^2$
checks against the momentum sum. The pair rapidity $Y=-2.595$ and the decay
rapidity $y^{*}=\operatorname{artanh}(\beta\cos\hat\theta)=1.157$ give
$y_b=Y+y^{*}=-1.438$ and $y_{\bar b}=Y-y^{*}=-3.752$, in agreement with the
table. Two quarks of identical \pt will have completely different fates. The
mother, the $b$, is born at $\eta=-1.44$: inside the barrel, where the
tracker, the electromagnetic and the hadron calorimeters will all see her.
The $\bar b$ at $\eta=-3.75$ heads into the forward hadron calorimeter, where
there are no tracks and therefore, much later in the journey, no $b$ tag.
Keep an eye on both.

\begin{figure}[htbp]\centering
\includegraphics[width=\textwidth]{ME/journey_ME.pdf}
\caption[The hard process of our event in the $(\eta,\phi)$ plane]{The hard
process of our event in the $(\eta,\phi)$ plane, drawn by
\file{example/ME/plot_journey_ME.py} from \file{example/ME/event_ME.lhe}.
The layout is that of the CMSSW \code{genPartAnalyzer} of the later
chapters: marker size grows with $\log\pt$; the faint bands are the CMS tracker
region, the endcap calorimeters and the forward calorimeter (HF). The
incoming partons and their momentum fractions are noted at the two edges.
Nothing has showered, hadronised or reached a detector yet.}
\label{fig:ancestral_home:journey}
\end{figure}

The ledger opens with a single entry: one parton, its \pt, its mass. Every
later row will be compared with this one.

\begin{ledger}{The Ancestral Home}
  \ledgerrow{before this stage}{--}{--}{--}{--}
  \ledgerrow{$b$ quark in the hard-process record}{113.22}{4.80}{1}{--}
\end{ledger}

\section{Further reading}
\label{sec:ancestral_home:further_reading}
The factorisation theorem is proved, as far as it is proved, in
\cite{Collins:1989gx}; the textbook account of everything in this chapter,
from the parton model to heavy-quark production, is \cite{Ellis:1996mzs},
and the Particle Data Group's ``Cross-section formulae'' and ``QCD'' reviews
\cite{ParticleDataGroup:2024cfk} give the massive-quark Born formula and the
running coupling in the conventions used here. The leading-order
heavy-quark cross sections are due to \cite{Combridge:1978kx}, their
next-to-leading-order completion to \cite{Nason:1987xz}. NNPDF3.1 is
described in \cite{NNPDF:2017mvq}, the LHAPDF library in
\cite{Buckley:2014ana}, and the PDF4LHC recommendations for combining PDF
uncertainties in \cite{Butterworth:2015oua}. The \PYTHIA~8.3
guide~\cite{Bierlich:2022pfr}, chapters on hard processes and on couplings
and scales, is the reference for process~124 and the scale codes; the CP5
tune is documented in \cite{CMS:2019csb}. \MG and its NLO capabilities are
described in \cite{Alwall:2014hca}. On the numerical side, the adaptive
importance sampling that generators use is \cite{Lepage:1977sw}, and the
Sobol sequence of the scripts goes back to \cite{Sobol:1967} with the
direction numbers of \cite{Joe:2008}.

\begin{pitfalls}
\begin{itemize}[nosep]
  \item \textbf{Four numbers, one cross section.} $345.3$, $345.1$ and
    $344.3\pb$ are the same calculation with different statistical
    fluctuations and conventions. Do not ``fix'' one to match another, and
    do not read the $0.3\%$ spread as physics: the theoretical uncertainty
    is $20\%$.
  \item \textbf{$\pthat$ is not jet \pt.} The generator cut acts on the two
    Born quarks. After showering, a $113\GeV$ quark becomes a jet of any \pt
    from a few GeV up; a sample with $\pthat\ge100\GeV$ is not a sample of
    jets above $100\GeV$, and cannot be used for jets near that threshold.
  \item \textbf{NNLO PDF, LO matrix element.} The order of the PDF set and the
    order of the calculation are independent choices. CMS uses NNLO PDFs with
    LO \PYTHIA processes; that is a tune decision, not a claim of NNLO accuracy.
  \item \textbf{Scale multipliers multiply $Q^2$.} \PYTHIA's
    \code{renormMultFac = 4} doubles $\muR$, not quadruples it. The
    integrator's \code{--mur-factor} option multiplies $\muR$ itself.
  \item \textbf{Which $\alphas$?} The matrix element, each shower, the
    multi-parton interactions and the PDF set each have their own coupling.
    Check \code{SigmaProcess:alphaSvalue} and \code{alphaSorder}, not
    \code{Tune:pp}. That they agree at our event is luck of the scale, not a
    rule.
  \item \textbf{$\mb=4.8\GeV$} is \PYTHIA's default mass for the $b$ in the
    matrix element; \MG's is $4.7\GeV$. Neither is the pole mass nor the
    $\overline{\mathrm{MS}}$ mass in the Particle Data Group tables, and
    changing it changes the threshold, the scale and (near threshold) the
    cross section.
  \item \textbf{$gg$ is off here.} An inclusive $\bbbar$ sample is
    dominated by gluon fusion. Our number covers only $q\bar q$ annihilation,
    and the \MG recipe removes the gluon from the proton by hand.
  \item \textbf{\code{define p} in \MG also removes the $b$.} The
    four-flavour proton of the recipe is a different scheme from \PYTHIA's
    five-flavour default, worth $1.6\%$ here.
  \item \textbf{The \file{.lhe} file is not an LHE file.} It is \PYTHIA's
    hard-process listing, saved with that extension because it plays the role
    of the LHE record. Status codes ($-21$, $23$) are \PYTHIA's, not the LHE
    standard's ($-1$, $1$).
  \item \textbf{Units.} \PYTHIA prints mb, \MG and the scripts print pb, and
    the scripts also print nb. $1\,\mathrm{mb}=10^{6}\nb=10^{9}\pb$; the
    $3.451\times10^{-7}\,\mathrm{mb}$ of the log is $345.1\pb$.
  \item \textbf{Which Python?} A shell alias such as
    \code{alias python=/usr/bin/python3} silently overrides a virtual
    environment. Run the scripts with the explicit interpreter path when in
    doubt; the standalone script prints the table provenance on its first
    line so you can see what was loaded.
\end{itemize}
\end{pitfalls}

\begin{exercises}
\item \textbf{Run it.} Run \file{standalone_qqbar_bbbar_xsec.py} with no
  options and confirm $345.320\pb$. Run it with \code{--repeats 8} and
  compare the error with \cref{tab:ancestral_home:three_numbers}. The
  \CMSSW job accepted $1000$ unweighted events and quoted a $1.8\%$
  statistical error; how many events would reach $0.1\%$? Why is that
  precision pointless for a leading-order calculation?
\item \textbf{The angular integral by hand.} Integrate
  \cref{eq:ancestral_home:dsigma} over $\cos\theta$ from $-1$ to $1$ and
  obtain \cref{eq:ancestral_home:sigmahat}. Evaluate both at our event
  ($\shat=156\,895\GeV^2$, $\alphas=0.1143$, $\cos\hat\theta=-0.820$) and
  confirm $\hat\sigma=30.2\pb$ and
  $\mathrm{d}\hat\sigma/\mathrm{d}\cos\theta=18.9\pb$. Then integrate only
  over $|\cos\theta|\le c_{\max}$ for $\pthat^{\min}=100\GeV$ and confirm
  that $81\%$ of $\hat\sigma$ survives the cut at this $\shat$.
\item \textbf{The Jacobian, once more.} Repeat the derivation of
  \cref{eq:ancestral_home:jacobian} for the variables $(\tau,Y)$ instead of
  $(\ln\tau,Y)$, and for $(\ln x_1,\ln x_2)$. Which of the three choices
  gives an integrand that is flattest in $u_1$ for a PDF behaving as
  $xf\sim x^{-0.3}$ at small $x$? This is why the script samples $\ln\tau$.
\item \textbf{Kinematics from the record.} Using only the four-vectors in
  \file{event_ME.lhe} and the beam energy, compute $x_1$, $x_2$, $\shat$, $Y$,
  $\pthat$, $\cos\hat\theta$, $y^{*}$ and the rapidities and pseudorapidities
  of both quarks. Compare with \cref{tab:ancestral_home:kinematics} and with
  the printout of \file{plot_journey_ME.py}. Why do $y$ and $\eta$ differ by
  only $0.001$ here, and when would they differ by much more?
\item \textbf{The \pt spectrum.} Reproduce \cref{tab:ancestral_home:ptscan}
  with \code{--ptmin 50}, \code{150}, \code{200} and \code{300}. Between
  $100$ and $300\GeV$ the cross section falls as $(\pthat^{\min})^{-n}$;
  determine $n$ in each interval. Dimensional analysis of
  \cref{eq:ancestral_home:sigmahat} suggests $n=4$; explain what the PDFs
  and the running coupling do to it, and why $n$ is smaller at low \pt.
\item \textbf{Where does the $\muF$ dependence come from?} Using the
  embedded PDF table (\code{TablePDF.xf} in the standalone script), evaluate
  $x f_{\bar u}(x,Q^2)$ at $x=2.17\times10^{-3}$ and $x f_u$ at $x=0.390$
  for $Q^2=12\,842/4$, $12\,842$ and $4\times12\,842\GeV^2$. One rises with
  $Q^2$ and one falls; by how much each? Connect this to the small size and
  the sign of the $\muF$ rows in \cref{tab:ancestral_home:variations}, and
  to the much larger $\muF$ dependence at $\pthat\ge30\GeV$, where
  $x_1$ was ten times smaller.
\item \textbf{Scale variation with one line.} In \code{born\_weights} the
  line \code{q2 = pt2 + mb2} sets both scales. Change it to
  \code{4*(pt2 + mb2)} and to \code{0.25*(pt2 + mb2)} and reproduce the two
  ``both scales'' rows of \cref{tab:ancestral_home:variations}. Why can the
  standalone script not reproduce the single-scale rows without a larger
  change? (Look at what \code{TablePDF.xf} and \code{TableAlphaS} are called
  with.)
\item \textbf{Match \MG.} Run the full script with \code{--nquark 4
  --mb 4.7} and then vary the scale code and multipliers until you
  reproduce $344.3\pb$ within \MG's error. Read \MG's run card
  documentation for \code{dynamical\_scale\_choice} to see which scale it
  actually used; is your match the right one or an accident?
\item \textbf{The incoming $b$.} Zero the fifth flavour in
  \code{born\_weights} (or run the full script with \code{--nquark 4}) and
  confirm the $1.6\%$ change. Draw the additional Feynman diagrams that
  $b\bar b\to b\bar b$ has and $u\bar u\to b\bar b$ does not. Why does the
  same-flavour case not matter much here, and when would it?
\item \textbf{The gluon channel (needs \CMSSW).} In \file{step1_GEN_cfg.py}
  switch the gluon channel on, \code{HardQCD:gg2bbbar = on}, run $1000$
  events and read the per-process cross sections from the statistics table
  in the log. What fraction of the inclusive $\bbbar$ rate at
  $\pthat\ge100\GeV$ is $q\bar q$ annihilation? Why is the $gg$ matrix
  element not a one-line formula?
\item \textbf{Film your own event.} Run the animation script with
  \code{--event} pointing at a hard-process record from your own \CMSSW log
  and \code{--scene collision}. Predict, before you watch, whether the pair
  in your event slides forward or backward along the tube, and by how much.
\item \textbf{Two fates.} The $\bar b$ of our event is at $\eta=-3.75$.
  Which CMS subdetectors will see it, and which will not? Will it become a
  jet in \cref{ch:the_inn}, and can it ever be $b$-tagged? Keep your
  answer; \cref{ch:registry_office} will check it.
\item \textbf{The learning mode by hand.} Run the standalone script with
  \code{-v} and reproduce section~[2] of its log with a pocket calculator:
  from $x_1$, $x_2$ and $\cos\theta^*$ of \file{event_ME.lhe} compute
  $(u_1,u_2,u_3)$, the three Jacobian factors, the ten PDF products and the
  five weights $w_q$. Where does the $650.7\pb$ come from, and why is it
  twice the mean weight although the event sits on the falling slope of
  \cref{fig:ancestral_home:factorisation_plot}?
\item \textbf{Sobol against random.} Replace the Sobol generator in
  \code{integrate} by NumPy's \code{default\_rng(seed)} and its
  \code{random((N,3))} method, and rerun with \code{--repeats 8}. Compare the error with
  \cref{tab:ancestral_home:three_numbers} and with panel~(c) of
  \cref{fig:ancestral_home:convergence_plot}. How many pseudo-random points
  would you need to reach the Sobol precision of $0.001\pb$, and why does the
  Sobol advantage shrink when the integrand has sharp edges, as it would with
  an upper cut \code{--ptmax} close to \code{--ptmin}?
\end{exercises}

\subsection*{Open problems in the context of Phase-2 LHC}
\label{sec:ancestral_home:open_problems}

The High-Luminosity LHC will collide protons at $\sqrt{s}=14\TeV$ with up to
200 simultaneous interactions per crossing and collect some
$3000\,\mathrm{fb}^{-1}$. The Standard Model programme that this makes
possible is surveyed in the HL-LHC Yellow Report~\cite{Azzi:2019yne}; the
expected gain in PDF precision in \cite{AbdulKhalek:2018rok}; and the strain
that the statistics put on event generation, in CPU time and in theoretical
accuracy, in the HSF generator report~\cite{HSFPhysicsEventGeneratorWG:2020gxw}.
Almost nothing at this first stop is affected by the pile-up; almost
everything is affected by the statistics. Here are six questions that are
open, or at least unsettled, at the time of writing, each with the state of
the art as I understand it. None of them needs more than the scripts of this
chapter to make a start, and each is the kind of study whose result the
generator authors and the CMS generator group would like to hear about.
\begin{enumerate}
  \item \textbf{The energy step.} Change \code{--ecm} to $14000$ and see
    how much of the $\bbbar$ rate at fixed $\pthat$ comes from the $3\%$
    extra energy and how much from the smaller $x$ it reaches. Then ask the
    harder question: the CP5 tune~\cite{CMS:2019csb} was fitted to $13\TeV$
    data, and its energy scaling of the multi-parton-interaction cut-off,
    \code{MultipartonInteractions:ecmPow}, is an extrapolation whose form
    goes back to the Monash tune~\cite{Skands:2014pea} and ultimately to the
    model of \cite{Sjostrand:1987su}. What early Run-4 measurement would pin
    it down, and how would the answer feed back into the underlying event of
    every jet in this book?
  \item \textbf{Which PDF set?} NNPDF3.1 is a 2017 fit~\cite{NNPDF:2017mvq}.
    Since then NNPDF4.0~\cite{NNPDF:2021njg}, CT18~\cite{Hou:2019efy} and
    MSHT20~\cite{Bailey:2020ooq} have appeared, and the PDF4LHC21
    combination~\cite{PDF4LHCWorkingGroup:2022cjn} averages them; the sets
    disagree most where our process is most sensitive, in the sea and in the
    heavy flavours, and NNPDF has since claimed evidence for an intrinsic
    charm component~\cite{Ball:2022qks}. Rerun the full script with
    \code{--pdf} set to each of these and compare the Born rate at
    $\pthat\ge100\GeV$ and its member-to-member spread. Which sets differ in
    the $b$ density, and does the 5FS incoming-$b$ term change by more than
    its own size between them? A Run-4 tune will be built on one of these;
    the choice is not yet made, and \cite{AbdulKhalek:2018rok} estimates how
    much HL-LHC data will narrow the spread.
  \item \textbf{Leading order is not enough, and NLO is expensive.} With
    $3000\,\mathrm{fb}^{-1}$ the statistical error of a $b$-jet measurement
    at $\pt\sim100\GeV$ will be far below the $20\%$ of
    \cref{eq:ancestral_home:envelope}, so production samples will have to
    be NLO~\cite{Nason:1987xz,Alwall:2014hca}, and for the cross section
    itself NNLO exists~\cite{Catani:2020kkl}. NLO event samples matched to a
    shower with the MC@NLO method contain negative-weight events, each of
    which costs a full detector simulation and cancels against a positive
    one; the HSF report~\cite{HSFPhysicsEventGeneratorWG:2020gxw} quantifies
    the CPU cost. Remedies exist on the generator side, the MC@NLO-$\Delta$
    prescription~\cite{Frederix:2020trv}, or the POWHEG method, which is
    positive by construction~\cite{Alioli:2010xd}, and on the sample side,
    positive resampling~\cite{Andersen:2020sjs}, cell
    resampling~\cite{Andersen:2021mvw} and neural
    resampling~\cite{Nachman:2020fff}. Generate $pp\to\bbbar$ at NLO in \MG
    at $14\TeV$ with our cut, measure the negative-weight fraction, estimate
    the simulation CPU it wastes per unit of statistical precision, and ask
    which of the remedies is usable in a CMS production.
  \item \textbf{Slicing the spectrum.} \Cref{tab:ancestral_home:ptscan}
    falls by four orders of magnitude between $30$ and $300\GeV$, with the
    local power $n$ rising from $3$ to $5$
    (\cref{fig:ancestral_home:flavours_plot}). A single sample cannot serve
    both ends, so HL-LHC productions are sliced in $\pthat$ or biased in
    \pt. Using \code{--ptmin} and \code{--ptmax}, design a set of $\pthat$
    slices for $\bbbar$ at $14\TeV$ whose events, after weighting, populate
    a jet \pt spectrum up to $2\TeV$ with a roughly constant statistical
    error, and estimate the total number of events to simulate. Compare with
    the biasing option \code{PhaseSpace:bias2SelectionPow} of
    \PYTHIA~\cite{Bierlich:2022pfr}, and with the generator-level
    efficiency considerations of \cite{HSFPhysicsEventGeneratorWG:2020gxw}.
  \item \textbf{The $b$ mass and the flavour scheme.} $\mb=4.8\GeV$ is a
    generator convention, and 4FS versus 5FS is a scheme choice that
    changes the answer at leading order by a few per cent
    (\cref{tab:ancestral_home:variations}). The two schemes are compared for
    $b$-initiated processes in \cite{Maltoni:2012pa} and matched into one in
    \cite{Forte:2015hba}; for $\bbbar$ production itself the FONLL
    predictions of \cite{Cacciari:2012ny} and the NNLO results of
    \cite{Catani:2020kkl} are the state of the art. Precision measurements
    with $b$ jets, such as $Z+b$ or $H\to\bbbar$, will need a consistent
    scheme across matrix element, PDF and shower. Compare 4FS and 5FS
    $\bbbar$ production in \MG at LO and NLO at $14\TeV$ and identify where
    in $\pthat$ they differ most.
  \item \textbf{The coupling.} The world average of $\alphas(m_Z)$ has
    moved and its uncertainty has shrunk since the tune was made; the 2024
    value is $0.1180\pm0.0009$~\cite{ParticleDataGroup:2024cfk}, and
    \cite{dEnterria:2022hzv} reviews how the next decade of measurements may
    halve the error. Propagate $\alphas(m_Z)=0.1180\pm0.0009$ through the
    full script with \code{--alphas-mz} and compare the shift with the PDF
    uncertainty and with \cref{eq:ancestral_home:scale_estimate}; then ask
    whether the four couplings of the CP5 tune (matrix element, two
    showers, MPI) should all move together, as the tune assumes.
\end{enumerate}
Later stops will add their own lists; the ones about the detector will be
longer.
